% PDF page 401
\source{gilbarg-trudinger:page-0401}

\[
\text{16.1. Hypersurfaces in } \mathbb R^{n+1}
\]

$\mathcal U$ of $\mathbb R^{n+1}$ and a function $\Phi \in C^2(\mathcal U)$, with $D\Phi \neq 0$ on $\mathcal U$, such that

\begin{equation}
\tag{16.2}
\mathcal S = \{x \in \mathcal U \mid \Phi(x)=0\}.
\end{equation}

For the applications of this chapter, $\mathcal S$ will always be the graph of a function
$u \in C^2(\Omega)$ where $\Omega$ is a domain in $\mathbb R^n$. In this case we can take $\mathcal U = \Omega \times \mathbb R$ and

\begin{equation}
\tag{16.3}
\Phi = x_{n+1} - u(x'), \qquad x' = (x_1, \ldots, x_n) \in \Omega.
\end{equation}

When $\mathcal S$ is given by (16.2) the normal $\nu$ to $\mathcal S$ (in the direction of increasing $\Phi$),
is given by

\[
\nu = \frac{D\Phi}{|D\Phi|}.
\]

For $g \in C^1(\mathcal U)$, the tangential gradient $\delta g$ of $g$ on $\mathcal S$ is defined by

\begin{equation}
\tag{16.4}
\delta g = Dg - (\nu \cdot Dg)\nu.
\end{equation}

For any point $y \in \mathcal S$, the tangential gradient $\delta g(y)$ is the projection of the gradient
$Dg(y)$ onto the tangent plane to $\mathcal S$ at $y$. Clearly we have

\begin{equation}
\tag{16.5}
\nu \cdot \delta g = 0
\end{equation}

and

\begin{equation}
\tag{16.6}
|\delta g|^2 = |Dg|^2 - |\nu \cdot Dg|^2,
\end{equation}

so that

\begin{equation}
\tag{16.7}
|\delta g| \le |Dg|
\end{equation}

and

\begin{equation}
\tag{16.8}
|\nu_{n+1}|\,|Dg| \le |\delta g| \qquad \text{if } D_{n+1}g \equiv 0.
\end{equation}

Furthermore, it is clear that $\delta g$ only depends on the values of $g$ on $\mathcal S$. For, suppose
$\bar g \in C^1(\mathcal U)$ satisfies $g=\bar g$ on $\mathcal S$. Then $D(g-\bar g)=k\nu$ for some constant $k$ and hence

\[
\delta(g-\bar g)=k\nu-k\nu=0.
\]

Next, by formula (14.102) in Section 14.6, we have

\[
\delta_i \nu_i = D_i\left[\frac{D_i\Phi}{|D\Phi|}\right] - \nu_i\nu_j D_j \nu_i
= -nH - \frac12 \nu_j D_j |\nu|^2
= -nH
\]

% PDF page 402
\source{gilbarg-trudinger:page-0402}

where $H$ denotes the mean curvature of $\mathcal{G}$ with respect to $v$. Hence we have the formula

\[
(16.9)\qquad H = - \frac{1}{n}\,\delta_i v_i.
\]

The following lemma provides an integration by parts formula for the differential operator $\delta$.

\medskip

\textbf{Lemma 16.1.} \emph{Letting $dA$ denote the area element in $\mathcal{G}$, we have}

\[
(16.10)\qquad \int_{\mathcal{G}} \delta g\, dA = -n \int_{\mathcal{G}} gHv\, dA
\]

\emph{for all $g \in C^1_0(\mathcal{U})$.}

\medskip

\textit{Proof.} We shall establish formula (16.10) for the case where $\mathcal{G}$ is the graph of a $C^2$ function. The general case then follows by means of a partition of unity. Accordingly, let us assume that $\Phi$ is given by (16.3) so that at points $(x,u(x)) \in \mathcal{G}$, we have

\[
v_i = - \frac{D_i u}{v}, \qquad i = 1,\dots,n,
\]

\begin{equation}
\tag{16.11}
v_{n+1} = \frac{1}{v},
\end{equation}

\[
H = \frac{1}{n} D_i \!\left[\frac{D_i u}{v}\right],
\]

and

\[
dA = v\, dx,
\]

where

\[
v = \sqrt{1 + |Du|^2}.
\]

Defining $\tilde g \in C^1(\mathcal{U})$ by

\[
\tilde g(x,x_{n+1}) = g(x,u(x)),
\]

we obviously have $\tilde g = g$ on $\mathcal{G}$ and hence $\delta \tilde g = \delta g$. Therefore, for $i \le n$,

\[
\int_{\mathcal{G}} \delta_i g\, dA = \int_{\mathcal{G}} \delta_i \tilde g\, dA
\]

\[
= \int_{\Omega} (D_i g - v_i v_j D_j g)\, v\, dx
\]

% PDF page 403
\source{gilbarg-trudinger:page-0403}

\setcounter{page}{391}

\[
= - \int_{\Omega} \tilde{g}\left\{D_i v - \sum_{j=1}^n D_j(v_i v_j)\right\}\,dx \qquad \text{by integration by parts,}
\]

\[
= - \int_{\Omega} \tilde{g}\left\{D_i v - \sum_{j=1}^n v_j D_j(v_i v) + n v H v_i\right\}\,dx \qquad \text{by (16.11),}
\]

\[
= - n \int_{\Omega} \tilde{g}Hv_i v\,dx - \int_{\Omega} \tilde{g}\left(\sum_{j=1}^n v_j D_i v_j + D_i v\right)\,dx
\]

\[
= - n \int_{\Omega} \tilde{g}Hv_i v\,dx \qquad \text{by (16.11).}
\]

For $i=n+1$, we have

\[
\int_{\mathcal G}\delta_{n+1}g\,dA=\int_{\mathcal G}\delta_{n+1}\tilde{g}\,dA
\]

\[
= - \int_{\Omega} v_{n+1}\sum_{j=1}^n v_jD_j\tilde{g}\,dx
\]

\[
= \int_{\Omega} \tilde{g}\sum_{j=1}^n D_jv_j\,dx
\]

\[
= - n \int_{\Omega}\tilde{g}H\,dx
\]

\[
= - n \int_{\mathcal G} gH v_{n+1}\,dA.\quad \square
\]

Note that the case $i=n+1$ in Lemma 16.1 is equivalent to the integral form of
the prescribed mean curvature equation.
For $g\in C^2(\mathcal U)$, the Laplacian (or Laplace-Beltrami operator) of $g$ on $\mathcal G$ is defined
by

\begin{equation}
\Delta g=\delta_i\delta_i g.
\tag{16.12}
\end{equation}

From the integration by parts formula (16.10) and (16.5) we have

\begin{equation}
\int_{\mathcal G}\varphi\,\Delta g\,dA=\int_{\mathcal G}g\,\Delta\varphi\,dA=\int_{\mathcal G}-\delta g\cdot\delta\varphi\,dA
\tag{16.13}
\end{equation}

for all $\varphi\in C_0^2(\mathcal U)$.


% PDF page 404
\source{gilbarg-trudinger:page-0404}

We proceed now to derive some important inequalities involving the operators
$\delta$ and $\Delta$ on $\mathcal{G}$. These inequalities will include useful extensions of the mean value
inequalities, Theorem 2.1, and the potential representation, Lemma 7.14, to
hypersurfaces in $\mathbb{R}^{n+1}$. If $y$ is a point on $\mathcal{G}$, $r = |x-y|$ and $\psi \in C^2(\mathbb{R})$ we have,
by calculation,

\[
\delta_i\psi(r) = (\delta_{ij} - v_i v_j)\frac{x_j-y_j}{r}\,\psi'(r),
\]

\begin{equation}
\Delta\psi(r) = \frac{n\psi'(r)}{r} + \left[\frac{\psi''(r)}{r^2} - \frac{\psi'(r)}{r^3}\right](r^2 - |v_i(x_i-y_i)|^2) + \frac{n\psi'(r)}{r} H v_i(x_i-y_i)
\tag{16.14}
\end{equation}
\[
= \frac{n\psi'(r)}{r} + \left[\psi'' - \frac{\psi'(r)}{r}\right]|\delta r|^2 + \frac{n\psi'(r)}{r} H v\cdot(x-y)
\]

since, by (16.6),

\[
|\delta r|^2 = 1 - \frac{|v\cdot(x-y)|^2}{r^2}.
\]

In particular, let $\chi$ be a non-negative, non-increasing function in $C^1(\mathbb{R})$ with
support in the interval $(-\infty, 1)$ and set

\[
\psi(r) = \int_r^\infty \tau\chi(\tau/\rho)\, d\tau
\]

where $0 < \rho < R$ and the ball $B_R(y) \subset \mathcal{U}$. We then have by (16.14)

\begin{equation}
\Delta\psi(r) = -\{n\chi(r/\rho) + r\chi'(r/\rho)|\delta r|^2/\rho + n\chi(r/\rho)Hv\cdot(x-y)\}
\tag{16.15}
\end{equation}
\[
= \rho^{n+1}D_\rho[\rho^{-n}\chi(r/\rho)] + r\chi'(r/\rho)(1-|\delta r|^2)/\rho - n\chi(r/\rho)Hv\cdot(x-y)
\]
\[
\leq \rho^{n+1}D_\rho[\rho^{-n}\chi(r/\rho)] - n\chi(r/\rho)Hv\cdot(x-y).
\]

In the special case where $\mathcal{G}$ is a minimal surface, that is $H \equiv 0$, inequality (16.15)
reduces to

\begin{equation}
\Delta\psi(r) \leq \rho^{n+1}D_\rho[\rho^{-n}\chi(r/\rho)].
\tag{16.16}
\end{equation}

The relations (16.15), (16.16) are fundamental to our treatment of interior esti-
mates. We illustrate their application by first considering the minimal surface case,
$H \equiv 0$. Let $g$ be a non-negative function in $L^1(\mathcal{G})$ and suppose that

\begin{equation}
\int_{\mathcal{G}} g\,\Delta\varphi\, dA \geq 0
\tag{16.17}
\end{equation}

% PDF page 405
\source{gilbarg-trudinger:page-0405}


for all non-negative $\phi \in C_0^2(\mathcal{U})$. By choosing $\phi=\psi$ in (16.17), we obtain immediately from (16.16) the inequality

\[
D_\rho\!\left[\frac{1}{\rho^n}\int_{\mathcal{S}} g\chi(r/\rho)\,dA\right]\geq 0,
\]

that is, the function $I_{\chi}$ given by

\begin{equation}
I_{\chi}(\rho)=\frac{1}{\omega_n\rho^n}\int_{\mathcal{S}} g\chi(r/\rho)\,dA
\tag{16.18}
\end{equation}

is non-decreasing in $\rho$. Letting $\chi$ approximate the characteristic function of the interval $(-\infty,1)$, in an appropriate fashion, we obtain that the function $I$ given by

\begin{equation}
I(\rho)=\frac{1}{\omega_n\rho^n}\int_{\mathcal{S}\cap B_{\rho}(y)} g\,dA
\tag{16.19}
\end{equation}

is also non-decreasing. Since

\begin{equation}
\lim_{\rho\to 0} I(\rho)=g(y)
\tag{16.20}
\end{equation}

for almost all (with respect to $dA$) $y\in\mathcal{S}$, we conclude the mean value inequality

\begin{equation}
g(y)\leq \frac{1}{\omega_n R^n}\int_{\mathcal{S}\cap B_R(y)} g\,dA
\tag{16.21}
\end{equation}

for almost all $y\in\mathcal{S}$ with $B_R(y)\subset \mathcal{U}$. Let us call a function $g\in C^2(\mathcal{U})$ subharmonic (harmonic) on $\mathcal{S}$ if $\Delta g\geq 0\; (=0)$ on $\mathcal{S}$. Using (16.13), we thus have

\textbf{Lemma 16.2.} \emph{Let $g$ be a non-negative, subharmonic function on a $C^2$ minimal hypersurface $\mathcal{S}$. Then for any point $y\in\mathcal{S}$ and ball $B_R(y)\subset \mathcal{U}$, the inequality (16.21) is valid.}

When $\mathcal{S}$ is a hyperplane, inequality (16.21) reduces to the mean value inequality for non-negative subharmonic functions in Euclidean space $\mathbb{R}^n$. Note however that in this case we do not need to assume that $g$ is non-negative. When $g$ is a positive constant, we obtain from (16.21) the estimate

\begin{equation}
A(\mathcal{S}\cap B_R(y))\geq \omega_n R^n
\tag{16.22}
\end{equation}

where $A(\mathcal{S}\cap B_R(y))$ denotes the area of $\mathcal{S}\cap B_R(y)$. Henceforth we shall write


% PDF page 406
\source{gilbarg-trudinger:page-0406}
$\Scal_R(y)=\Scal\cap B_R(y)$ and generally abbreviate $\Scal_R(y)=\Scal_R$ when there is no ambiguity.

We turn our attention now to the case of a general hypersurface \(\Scal\) and \(C^1(\Ucal)\) function \(g\). Although the procedure adopted above in the minimal surface case can be generalized, we shall instead proceed somewhat differently, thereby giving an alternative proof of Lemma 16.2. Let \(y\) be a point on \(\Scal\) and suppose that the ball \(B_R(y)\subset\Ucal\). Let \(\chi\) be a non-negative, non-increasing function in \(C^1[0,\infty)\) with support in the interval \([0,R]\). Defining \(\psi\in C^2(\Ucal)\) by

\[
\psi(r)=\int_r^\infty \tau\chi(\tau)\,d\tau,
\]

we have, by (16.15),

\begin{equation}\tag{16.23}
\Delta\psi(r)= -\{n\chi(r)+r\chi'(r)|\delta r|^2+n\chi(r)Hv\cdot (x-y)\}
=-(n\chi(r)+r\chi'(r))+r\chi'(r)(1-|\delta r|^2)-n\chi(r)Hv\cdot(x-y).
\end{equation}

Therefore, substituting into (16.13), we obtain

\begin{equation}\tag{16.24}
\int_{\Scal}\bigl(n\chi(r)+r\chi'(r)\bigr)g\,dA-\int_{\Scal}r\chi'(r)(1-|\delta r|^2)g\,dA
=-\,n\int_{\Scal}\chi(r)gHv\cdot(x-y)\,dA+\int_{\Scal}\delta\psi\cdot\delta g\,dA.
\end{equation}

If we further restrict \(\chi\) so that

\[
\chi(r)=\varepsilon^{-n}-R^{-n}\qquad\text{for }r<\varepsilon,
\]

where \(0<\varepsilon<R\), we can write the above relation as

\begin{equation}\tag{16.25}
n(\varepsilon^{-n}-R^{-n})\int_{\Scal_\varepsilon}g\,dA+\int_{\Scal_R-\Scal_\varepsilon}\bigl(n\chi(r)+r\chi'(r)\bigr)g\,dA
-\int_{\Scal_R-\Scal_\varepsilon}r\chi'(r)(1-|\delta r|^2)g\,dA
=-\,n\int_{\Scal_R}\chi(r)gHv\cdot(x-y)\,dA+\int_{\Scal_R}\delta\psi\cdot\delta g\,dA.
\end{equation}

% PDF page 407
\source{gilbarg-trudinger:page-0407}

The form of (16.25) suggests that we should choose $\chi$ so that $n\chi(r)+r\chi'(r)=\mathrm{con}$-
stant in the interval $[\varepsilon, R)$. Accordingly, let us define a function $\chi_\varepsilon$ by

\[
\chi_\varepsilon(r)=
\begin{cases}
\varepsilon^{-n}-R^{-n} & \text{for } 0\leq r<\varepsilon,\\
r^{-n}-R^{-n} & \text{for } \varepsilon\leq r<R,\\
0 & \text{for } r\geq R.
\end{cases}
\]

We cannot immediately replace $\chi$ by $\chi_\varepsilon$ in (16.25). However, this can be accom-
plished by replacing $\chi$ by a sequence $\{\chi_m\}$ of non-negative, non-increasing functions
in $C^1[0,\infty)$, with support in $[0,R)$ and uniformly bounded derivatives. By further
requiring that $\{\chi_m\}$ converges uniformly to $\chi_\varepsilon$ and that the sequence of derivatives
$\{\chi'_m\}$ converges pointwise to the function

\[
\chi'_\varepsilon(r)=
\begin{cases}
0 & \text{for } 0\leq r<\varepsilon,\\
-nr^{-(n+1)} & \text{for } \varepsilon\leq r<R,\\
0 & \text{for } r\geq R,
\end{cases}
\]

we can conclude, from (16.25),

\[
(\varepsilon^{-n}-R^{-n})\int_{\Ss_\varepsilon} g\,dA+\int_{\Ss_R-\Ss_\varepsilon} gr^{-n}(1-|\delta r|^2)\,dA
\]
\begin{equation}
=R^{-n}\int_{\Ss_R-\Ss_\varepsilon} g\,dA-(\varepsilon^{-n}-R^{-n})\int_{\Ss_\varepsilon} g\Hv\cdot(x-y)\,dA
\notag\\
-\int_{\Ss_R-\Ss_\varepsilon} g(r^{-n}-R^{-n})\Hv\cdot(x-y)\,dA+\frac{1}{n}\int_{\Ss_R}\delta\psi_\varepsilon(r)\cdot\delta g\,dA,
\tag{16.26}
\end{equation}

where

\[
\psi_\varepsilon(r)=\int_r^R t\chi_\varepsilon(t)\,dt.
\]

Letting $\varepsilon$ tend to zero, we thus have

\[
g(y)+\frac{1}{\omega_n}\int_{\Ss_R}\frac{(1-|\delta r|^2)}{r^n}\,g\,dA
\]
\begin{equation}
=\frac{1}{\omega_n R^n}\int_{\Ss_R} g\,dA-\frac{1}{\omega_n}\int_{\Ss_R} g(r^{-n}-R^{-n})\Hv\cdot(x-y)\,dA
\notag\\
+\frac{1}{n\omega_n}\int_{\Ss_R}\delta\psi(r)\cdot\delta g\,dA,
\tag{16.27}
\end{equation}

% PDF page 408
\source{gilbarg-trudinger:page-0408}
where

\[
\psi(r)=\int_r^R \tau\bigl(\tau^{-n}-R^{-n}\bigr)\,d\tau.
\]

Lemma 16.2 is now an immediate consequence of the identities (16.27) and (16.13).
Using the inequality

\begin{equation}
|Hv\cdot (x-y)|\le r^{-2}|v\cdot (x-y)|^2+\frac14 H^2r^2
=1-|\delta r|^2+\frac14 H^2r^2
\tag{16.28}
\end{equation}

we can deduce, from (16.27), the estimate

\begin{equation}
g(y)\le \frac{1}{\omega_n R^n}\int_{\mathcal{G}_R} g|\delta r|^2\,dA
+\frac{1}{4\omega_n}\int_{\mathcal{G}_R} gH^2r^2(r^{-n}-R^{-n})\,dA
-\frac{1}{n\omega_n}\int_{\mathcal{G}_R}(r^{-n}-R^{-n})(x-y)\cdot \delta g\,dA.
\tag{16.29}
\end{equation}

Thus we have proved the following generalization of Lemma 16.2.

\textbf{Lemma 16.3.} \textit{Let $g$ be a non-negative function in $C^1(\mathcal{U})$. Then for any point $y\in \mathcal{G}$
and ball $B_R(y)\subset \mathcal{U}$, the inequality (16.29) is valid.}

The derivation of the interior gradient bound in the following section will be
based upon Lemma 16.3. Note that for $g\in C^2(\mathcal{U})$ we can express the last term in
(16.29) as

\[
-\frac{1}{n\omega_n}\int_{\mathcal{G}_R}\psi(r)\,\Delta g\,dA.
\]

Further consequences of inequalities (16.26) and (16.29) will be required for the
treatment of equations in two variables later in this chapter. In particular, if we set
$n=2$ in inequality (16.29), we obtain

\begin{equation}
g(y)\le \frac{1}{\pi R^2}\int_{\mathcal{G}_R} g\left(1+\frac14 H^2R^2\right)\,dA
+\frac{1}{2\pi}\int_{\mathcal{G}_R} r(r^{-2}-R^{-2})|\delta g|\,dA.
\tag{16.30}
\end{equation}

Setting $g\equiv 1$ in (16.30), we have the estimate

\begin{equation}
1\le \frac{A(\mathcal{G}_R)}{\pi R^2}+\frac{1}{4\pi}\int_{\mathcal{G}_R} H^2\,dA.
\tag{16.31}
\end{equation}

% PDF page 409
\source{gilbarg-trudinger:page-0409}

where $A(\mathcal{S}_R)$ denotes the area of $\mathcal{S}_R$. Consequently, if $\mathcal{S}$ is a compact hypersurface
in $\mathbb{R}^3$ (or more generally if $\mathcal{U}=\mathbb{R}^3$ and $A(\mathcal{S}_R)=o(R^2)$ as $R\to\infty$), we have

\begin{equation}
\int_{\mathcal{S}} H^2\,dA \geq 4\pi.
\tag{16.32}
\end{equation}

Furthermore, one can show that equality holds in (16.32) if and only if $\mathcal{S}$ is a
sphere. [TR 9].
Next, defining $I$ by (16.19), we obtain from (16.26) and (16.28)

\begin{align*}
\omega_n I(\epsilon)&=\epsilon^{-n}\int_{\mathcal{S}_\epsilon} g\,dA\\
&\leq R^{-n}\int_{\mathcal{S}_R} g|\delta r|^2\,dA+\epsilon^{1-n}\int_{\mathcal{S}_\epsilon} g|H|\,dA\\
&\quad+\frac{1}{4}\int_{\mathcal{S}_R} gH^2r^2(r^{-n}-R^{-n})\,dA+\frac{1}{n}\int_{\mathcal{S}_R} \delta\psi_\epsilon(r)\cdot\delta g\,dA.
\end{align*}

Hence, using Young's inequality (7.6), we obtain the estimate

\begin{align}
\sup_{(0,R)} I(\rho)&\leq \frac{1}{\omega_n R^n}\int_{\mathcal{S}_R}\big[g|n\delta r|^2+|HR|^n+nH^2r^2R^n(r^{-n}-R^{-n})\big]\,dA\\
&\quad+\frac{1}{\omega_n}\int_{\mathcal{S}_R} r(r^{-n}-R^{-n})|\delta g|\,dA.
\tag{16.33}
\end{align}

Note that the last term in (16.33) can be replaced by

\begin{equation*}
\frac{1}{\omega_n}\int_{\mathcal{S}_R}\psi(r)(-\Delta g)^+\,dA.
\end{equation*}

when $g\in C^2(\mathcal{U})$. Specializing inequality (16.33) to the case $n=2$, $g=1$, we obtain
the estimate

\begin{equation}
\sup_{(0,R)} \frac{A(\mathcal{S}_\rho)}{\rho^2}\leq \frac{2A(\mathcal{S}_R)}{R^2}+3\int_{\mathcal{S}_R} H^2\,dA.
\tag{16.34}
\end{equation}

We next determine an estimate for the quantity

\[
J(\rho)=D_\rho\left[\int_{\mathcal{S}_\rho} g|\delta r|^2\,dA\right].
\]

% PDF page 410
\source{gilbarg-trudinger:page-0410}
Choosing $\chi$ as in the proof of Lemma 16.2 and using (16.13) and (16.5) we have

\[
D_{\rho}\left[\int_{\mathcal S}\chi(r/\rho)|\delta r|^{2}\,dA\right]
\le \frac{n}{\rho}\int_{\mathcal S}\chi(r/\rho)\bigl[1+H\mathbf v\cdot(x-y)\bigr]g\,dA
\]

\[
-\frac{1}{\rho}\int_{\mathcal S}\delta\psi(r)\cdot\delta g\,dA.
\]

Hence as $\chi$ approaches the characteristic function of the interval $(-\infty,1)$, we see that $J(\rho)$ is well defined for all $\rho\in (0,R)$ and moreover that

\begin{equation}
J(\rho)\le \frac{1}{\rho}\int_{\mathcal S_{\rho}}\bigl[n|g|(1+|H|r)+r|\delta g|\bigr]\,dA
\tag{16.35}
\end{equation}

\[
\le \frac{n}{\rho}\int_{\mathcal S_{\rho}}|g|\,dA+\int_{\mathcal S_{\rho}}(n|Hg|+|\delta g|)\,dA.
\]

In particular, for $g=1$, we have

\[
D_{\rho}\left[\int_{\mathcal S_{\rho}}|\delta r|^{2}\,dA\right]
\le n\left[\frac{A(\mathcal G_{\rho})}{\rho}+\int_{\mathcal S_{\rho}}|H|\,dA\right]
\]

\begin{equation}
\le 10\rho\left[\frac{A(\mathcal G_{\rho})}{R^{2}}+\int_{\mathcal S_{R}}H^{2}\,dA\right]
\tag{16.36}
\end{equation}

by (16.34),

if $n=2$.

The potential type relations (16.27), (16.29) and (16.30) are somewhat analogous to Lemmas 7.14 and 7.16 in Chapter 7. Indeed, we shall now use (16.30) to derive an analogue of the Morrey estimate, Theorem 7.19, for two dimensional surfaces. The following lemma will be applied in Section 16.5 to conclude a Hölder estimate for generalized quasiconformal mappings.

\textbf{Lemma 16.4.} \emph{Let $g\in C^{1}(\mathcal U)$, $n=2$ and suppose there exist constants $K>0$ and $\beta\in(0,1)$ such that}

\begin{equation}
\int_{\mathcal S_{\rho}(y)}|\delta g|\,dA\le K\rho(\rho/R)^{\beta}
\tag{16.37}
\end{equation}

\emph{for all $\bar y\in \mathcal G_{R/4}(y)$ and all $\rho\le R/4$. Then}

\begin{equation}
\sup_{x\in \mathcal S_{\rho}(y)}|g(x)-g(y)|\le CK(\rho/R)^{\beta}\left\{\frac{A(\mathcal G_{R})}{R^{2}}+\int_{\mathcal S_{R}}H^{2}\,dA\right\}
\tag{16.38}
\end{equation}

% PDF page 411
\source{gilbarg-trudinger:page-0411}

where $C$ is a constant, and where $\mathcal{S}^*_\rho(y)$ denotes the component of $\mathcal{S}_\rho(y)$ which
contains $y$.

\textit{Proof.} We commence by writing (16.30) in the form

\begin{equation}
g(y)\le \frac{1}{\pi}\left\{\int_{\mathcal{S}_R} g\left(R^{-2}+H^2/4\right)\,dA
+ \int_0^R \rho^{-2}\int_{\mathcal{S}_\rho} |\delta g|\,dA\,d\rho\right\},
\tag{16.39}
\end{equation}

and defining, for $\rho\leq R/4$,

\[
g_1=\sup_{\mathcal{S}^*_\rho(y)} g,\qquad g_0=\inf_{\mathcal{S}^*_\rho(y)} g.
\]

If

\[
g_1-g_0\le 6\beta^{-1}K(\rho/R)^\beta,
\]

then Lemma 16.4 is established with $C=6\beta^{-1}$ by (16.34). If

\[
g_1-g_0>6\beta^{-1}K(\rho/R)^\beta,
\]

then we let $N$ be the largest integer such that

\[
N\le \beta(g_1-g_0)/6K(\rho/R)^\beta,
\]

and we subdivide the interval $[g_0,g_1]$ into $N$ pairwise disjoint intervals $I_1,I_2,\ldots,I_N$
of length $\ge 6\beta^{-1}K(\rho/R)^\beta$. For each $j=1,\ldots,N$, we then let $\psi_j$ be a non-negative
$C^1(\mathbb{R})$ function with support contained in $I_j$, max $\psi_j=1$ and max $\psi_j'\le \beta/2K(\rho/R)^\beta$.
(It is clear that such a function exists, because length $I_j\ge 6\beta^{-1}K(\rho/R)^\beta$.) Since $\mathcal{S}^*_\rho(y)$
is connected, we know that for each $j=1,\ldots,N$, there is a point $x^{(j)}\in \mathcal{S}^*_\rho(y)$ such
that $\psi_j[g(x^{(j)})]=1$. Then using (16.39) with $x^{(j)}$ in place of $y$, $\rho$ in place of $R$ and
$\psi_j\circ g$ in place of $g$, we obtain

\[
1\le \frac{1}{\pi}\int_{\mathcal{S}_{\sigma(x^{(j)})}} \psi_j(g)\left(\rho^{-2}+H^2/4\right)\,dA
+ \int_0^\rho \sigma^{-2}\int_{\mathcal{S}_{\sigma(x^{(j)})}} |\delta\psi_j\circ g|\,dA\,d\sigma,
\]

for all $\rho\le R/4$. Consequently, using (16.37), we obtain

\[
\int_0^\rho \sigma^{-2}\int_{\mathcal{S}_{\sigma(x^{(j)})}} |\delta\psi_j\circ g|\,dA\,d\sigma
\le \beta R^\beta(2K\rho^\beta)^{-1}\int_0^\rho \sigma^{-2}\int_{\mathcal{S}_{\sigma(x^{(j)})}} |\delta g|\,dA\,d\sigma
\]
\[
\le \beta R^\beta(2K\rho^\beta)^{-1}KR^{-\beta}\int_0^\rho \sigma^{\beta-1}\,d\sigma
\]
\[
=\frac{1}{2}.
\]

% PDF page 412
\source{gilbarg-trudinger:page-0412}

Combining the last two inequalities then gives

\[
1 \leq \frac{1}{\pi} \int_{\Srho(x^{(j)})} \psi_j(g)(\rho^{-2} + H^2/4)\,dA + \frac{1}{2}.
\]

so that

\[
1 \leq \frac{2}{\pi} \int_{\Srho(x^{(j)})} \psi_j(g)(\rho^{-2} + H^2/4)\,dA
\]

\[
\leq \frac{2}{\pi} \int_{\mathcal{S}_{2,\mu}(y)} \psi_j(g)(\rho^{-2} + H^2/4)\,dA.
\]

Summing over $j=1,\ldots,N$, and noting that $\sum \psi_j \leq 1$, we then deduce

\[
N \leq \frac{2}{\pi} \int_{\mathcal{S}_{2,\mu}(y)} (\rho^{-2} + H^2)\,dA
\]

\[
\leq C\left\{\frac{A(\SR)}{R^2} + \int_{\SR} H^2\,dA\right\} \quad \text{by (16.34)}
\]

and hence the estimate (16.38) follows.  $\square$

\emph{Remark.} If the function $g$ has compact support, we obtain, by letting $R$ approach
infinity in (16.27) and (16.29), the relations

\[
g(y) + \frac{1}{\wn} \int_{\mathcal{S}} \frac{(1-\lvert \delta r\rvert^2)}{r^n}\,g\,dA
= \frac{1}{\wn} \int_{\mathcal{S}} \frac{H\,v\cdot(x-y)}{r^n}\,g\,dA
- \frac{1}{n\wn} \int_{\mathcal{S}} \frac{(x-y)\cdot \delta g}{r^n}\,dA.
\]

\[
(16.40)
\]

\[
g(y) \leq \frac{1}{4\wn} \int_{\mathcal{S}} gH\,r^{2-n}\,dA
- \frac{1}{n\wn} \int_{\mathcal{S}} \frac{(x-y)\cdot \delta g}{r^n}\,dA.
\]

These inequalities can be used to establish imbedding theorems of the Sobolev
type (Problems 16.1, 16.2). We remark that the Sobolev inequality, Theorem 7.10,
has been extended to minimal hypersurfaces in $\Rn$ by Miranda [MD 1] and to
arbitrary submanifolds by Allard [AA], Michael and Simon [MS].

% PDF page 413
\source{gilbarg-trudinger:page-0413}

\textbf{16.2. Interior Gradient Bounds}

Let $\Omega$ be a domain in $\mathbb{R}^n$ and $u$ a function in $C^2(\Omega)$. If $\mathcal{G}$ denotes the graph of $u$
in $\mathbb{R}^{n+1}$, then the mean curvature (with respect to the upper normal) of $\mathcal{G}$, at the
point $x=(x',u(x'))\in\mathcal{G}$, is given by

\[
\tag{16.41}
H(x')=-\frac{1}{n}D_i\nu_i(x')
=\frac{1}{n}D_i\!\left[\frac{D_i u}{v}\right](x'),
\]

where $v=\sqrt{1+|Du|^2}$.

The object of this section is to establish a bound for $Du(x')$ in terms of $H$ and
$\operatorname{dist}(x',\partial\Omega)$. We commence by writing the relation (16.41) in the integral form

\[
\tag{16.42}
\int_{\Omega} (v\cdot D\varphi-nH\varphi)\,dx'=0
\]

for all $\varphi\in C_0^1(\Omega)$. Replacing $\varphi$ by $D_k\varphi$, $k=1,\ldots,n$, and integrating by parts, we
obtain

\[
\tag{16.43}
\int_{\Omega} (D_kv\cdot D\varphi-nD_kH\varphi)\,dx'=0
\]

for all $\varphi\in C_0^1(\Omega)$ (cf. equation (13.3)). In (16.43) and henceforth we assume that
$H\in C^1(\Omega)$. Now let us replace $\varphi$ in (16.43) by $v_k\varphi$ where $\varphi\in C_0^1(\Omega)$. Then we obtain
the equation

\[
\tag{16.44}
\int_{\Omega} D_k v_i D_i v_k\varphi\,dx'
+\int_{\Omega} (v_k D_k v_i D_i\varphi-nv_kD_kH\varphi)\,dx'=0
\]

for all $\varphi\in C_0^1(\Omega)$. By calculation we have

\[
v_kD_kv_i=-\frac{v_k}{v}(\delta_{ij}-v_iv_j)D_{jk}u
=-v(\delta_{ij}-v_iv_j)D_j\!\left(\frac{1}{v}\right)
=-v\,\delta_i v_{n+1}
\]

by (15.4), $i=1,\ldots,n$. Hence, by means of formula (14.104) we can write the above

% PDF page 414
\source{gilbarg-trudinger:page-0414}

relation as

\[
\tag{16.45}
\int_{\Omega} \mathcal{C}^{2}\varphi\,dx' - \int_{\Omega} v(\delta_i v_{n+1} D_i\varphi - n\delta_{n+1} H\varphi)\,dx' = 0,
\]

where

\[
\mathcal{C}^{2} = \sum_{i,j=1}^{n} D_i v_j D_j v_i = \sum_{i,j=1}^{n+1} (\delta_i v_j)^{2}
\]

Next, let us set $\mathcal{U} = \Omega \times \mathbb{R}$ and suppose that $\varphi \in C_0^1(\mathcal{U})$. By replacing $\varphi$ in (16.45)
by the function

\[
\tilde{\varphi}(x', x_{n+1}) = \varphi(x', u(x')),
\]

and using the relations

\[
\sum_{i=1}^{n+1} \delta_i v_{n+1}\delta_i\varphi = \sum_{i=1}^{n+1} \delta_i v_{n+1}\delta_i\tilde{\varphi} = \sum_{i=1}^{n} \delta_i v_{n+1} D_i\tilde{\varphi},
\]

we can therefore conclude from (16.45) the identity

\[
\tag{16.46}
\int_{\mathcal{G}} (\mathcal{C}^{2}v_{n+1}\varphi - \delta_i v_{n+1}\delta_i\varphi + n\delta_{n+1} H\varphi)\, dA = 0
\]

for all $\varphi \in C_0^1(\mathcal{U})$. Note that the functions $v$ and $H$ in (16.46) are independent of
$x_{n+1}$. Also, in (16.46) and in the remainder of this section, we follow the summation
convention that repeated indices indicate summation from $1$ to $n+1$. Defining the
function $w$ by

\[
\tag{16.47}
w = \log v = -\log v_{n+1},
\]

we obtain, by replacing $\varphi$ by $\varphi v$ in (16.46), the inequality

\[
\tag{16.48}
\int_{\mathcal{G}} (\delta w \cdot \delta \varphi + |\delta w|^{2}\varphi - nv \cdot D H\varphi)\, dA \leq 0
\]

for all non-negative $\varphi \in C_0^1(\mathcal{U})$. Inequality (16.48) is the weak form of the inequality

\[
\tag{16.49}
\Delta w \geq |\delta w|^{2} - nv \cdot D H.
\]

In particular, if $\mathcal{G}$ is a minimal surface, that is, the function $u$ satisfies the minimal
surface equation in $\Omega$, we see that $w$ is weakly subharmonic on $\mathcal{G}$ and hence Lemma
16.2 is applicable to $w$. In the general case, we can apply Lemma 16.3 to obtain, for

% PDF page 415
\source{gilbarg-trudinger:page-0415}

any point $y' \in \Omega,$

\begin{align*}
w(y') &\leq \frac{1}{\omega_n R^n} \int_{\mathcal{G}_{R(y)}} w\, dA + \frac{1}{4\omega_n} \int_{\mathcal{G}_{R(y)}} w H^2 r^{2}(r^{-n} - R^{-n})\, dA \tag{16.50}\\
&\qquad + \frac{1}{\omega_n} \int_{\mathcal{G}_{R(y)}} \psi(r)\, v \cdot DH\, dA
\end{align*}

provided $R < \operatorname{dist}(y', \partial\Omega)$, $y = (y', u(y'))$ and $\psi$ is given by (16.27). Setting

\begin{equation}\tag{16.51}
H_0 = \sup_{\Omega} |H|,\qquad
H_1 = \sup_{\Omega} (v \cdot DH)^+ \leq \sup_{\Omega} |DH|,
\end{equation}

we have from (16.50)

\begin{align*}
w(y') &\leq \frac{1}{\omega_n R^n} \int_{\mathcal{G}_R} w\, dA + \frac{H_0^2}{4\omega_n} \int_{\mathcal{G}_R} w r^2(r^{-n} - R^{-n})\, dA \\
&\qquad + \frac{H_1}{\omega_n} \int_{\mathcal{G}_R} \psi(r)\, dA \\
&= \frac{1}{\omega_n R^n} \int_{\mathcal{G}_R} w\, dA + \frac{n H_0^2}{4\omega_n} \int_0^R \rho^{-n-1} \int_{\mathcal{G}_\rho} w r^2\, dA\, d\rho \\
&\qquad + \frac{H_1}{\omega_n} \int_0^R \rho(\rho^{-n} - R^{-n}) A(\mathcal{G}_\rho)\, d\rho \qquad \text{(by Fubini's theorem)} \tag{16.52}\\
&\leq \frac{1}{\omega_n R^n} \int_{\mathcal{G}_R} w\, dA + \frac{n H_0^2}{4\omega_n} \int_0^R \rho^{1-n} \int_{\mathcal{G}_\rho} w\, dA\, d\rho \\
&\qquad + \frac{H_1}{\omega_n} \int_0^R \rho^{1-n} A(\mathcal{G}_\rho)\, d\rho.
\end{align*}

The estimation of $w$, and consequently $Du$, is thus reduced to the estimation of

\[
A(\mathcal{G}_\rho) \text{ and } \int_{\mathcal{G}_\rho} w\, dA \text{ for } 0 < \rho \leq R.
\]


% PDF page 416
\source{gilbarg-trudinger:page-0416}

\textit{Estimation of $A(\mathcal{G}_\rho)$.}

Let us assume that $3R < \operatorname{dist}(y', \partial\Omega)$ and also, without loss of generality, that
$y' = 0$ and $u(y') = 0$. For $\rho \leq R$, we define the function $u_\rho$ by

\[
u_\rho =
\begin{cases}
\rho & \text{for } u \geq \rho \\
u & \text{for } -\rho \leq u \leq \rho \\
-\rho & \text{for } u \leq -\rho
\end{cases}
\]

and substitute the test function

\[
\varphi = \eta u_\rho
\]

into the integral identity (16.42), where $\eta$ is a uniformly Lipschitz continuous function satisfying $\eta \equiv 1$ for $|x'| < \rho$, $\eta \equiv 0$ for $|x'| > 2\rho$ and $|D\eta| \leq 1/\rho$. Note that the identity (16.42) clearly holds for all $\varphi \in W_0^{1,1}(\Omega)$ and hence for all uniformly Lipschitz continuous $\varphi$ with support in $\Omega$. We obtain thus

\[
\int_{|x'|,\,|u|<\rho} \frac{|Du|^2}{v}\, dx' \leq \rho \int_{|x'|<2\rho} \left(|D\eta| + \eta |H| \eta \right) dx'.
\]

Consequently

\begin{align*}
A(\mathcal{G}_\rho) &= \int_{|x'|^2 + u^2 < \rho^2} v\, dx' \\
&\leq \int_{|x'|,\,|u|<\rho} v\, dx' \tag{16.53} \\
&\leq C(n)\left\{\rho^n + \rho \int_{|x'|<2\rho} |H|\eta\, dx'\right\} \\
&\leq C(n)\rho^n(1 + H_0 R).
\end{align*}

\textit{Estimation of $\int_{\mathcal{G}_\rho} w\, dA$.}

Let us now substitute

\[
\varphi = \eta w(u_\rho + \rho)
\]


% PDF page 417
\source{gilbarg-trudinger:page-0417}

\pagestyle{empty}



\setcounter{page}{405}

\noindent\textbf{16.2. Interior Gradient Bounds}\hfill 405

\bigskip

into (16.42) where $\eta$ is as above. We obtain thus

\[
\int_{|x'|,\,|u|<\rho} \frac{w|Du|^2}{v}\,\dxp
\le 2\rho \int_{|x'|<2\rho,\,u>-\rho} (w|D\eta|+\eta|Dw|+\eta|Hw\eta|)\,\dxp.
\]

In order to estimate $\int \eta|Dw|\,\dxp$, we replace $\varphi$ by $\varphi^2$ in inequality (16.48) so that

\[
\int_{\mathcal{G}} \varphi^2|w|^2\,dA
\le -2\int_{\mathcal{G}} \varphi\,\delta w\cdot\delta\varphi\,dA
- n\int_{\mathcal{G}} \varphi^2\,v\cdot DH\,dA
\]

for all $\varphi\in C_0^1(\Omega\times\mathbb{R})$. Using Cauchy's inequality (7.6), we obtain

\[
\int_{\mathcal{G}} \varphi^2|\delta w|^2\,dA
\le 4\int_{\mathcal{G}} |\delta\varphi|^2\,dA + nH_1 \int_{\mathcal{G}} \varphi^2\,dA.
\]

In particular, let us choose $\varphi$ such that

\[
\varphi=\eta\tau(x_{n+1}),
\]

where $0\le \tau\le 1$, $\tau\equiv 1$ in $\bigl(-\rho,\sup_{|x'|<2\rho}u\bigr)$, $\tau\equiv 0$ outside $\bigl(-2\rho,\rho+\sup_{|x'|<2\rho}u\bigr)$, $|\tau'|<2/\rho$
and $\eta$ is as above. We then have

\[
\int_{\mathcal{G}} \varphi^2|\delta w|^2\,dA \le (8\rho^{-2}+nH_1)\,A(\mathcal{G}\cap\supp\varphi).
\]

Using (16.8), we can then conclude that

\[
\int_{|x'|<2\rho,\,u>-\rho}\eta|Dw|\,\dxp
\le \int_{\mathcal{G}} \varphi|\delta w|\,dA
\]

\[
\le (8\rho^{-2}+nH_1)^{1/2}A(\mathcal{G}\cap\supp\varphi)
\quad\text{by Schwarz's inequality.}
\]

\[
\le (8+nH_1R^2)^{1/2}\rho^{-1}
\int_{|x'|<2\rho,\,u>-\rho} v\,\dxp.
\]

Since $w\le v$, we also have

\[
\int_{|x'|<2\rho,\,u>-\rho} w\,\dxp \le \int_{|x'|<2\rho,\,u>-\rho} v\,\dxp.
\]

It therefore only remains to estimate $\int v\,\dxp$, and this we accomplish by taking

\[
\varphi=\eta\max\{u+\zeta_\rho,0\}
\]


% PDF page 418
\source{gilbarg-trudinger:page-0418}

\pagestyle{empty}



\setcounter{page}{406}

\noindent\textbf{16. Equations of Mean Curvature Type}\hfill 406

\bigskip

in (16.42), where $\eta=1$ for $|x'|<2\rho$, $\eta=0$ for $|x'|>3\rho$ and $|D\eta|<1/\rho$. We obtain accordingly

\[
\int_{|x'|<2\rho,\,w>-\rho} v\,\dxp \le C(n)\rho^n(1+H_0R)\bigl(1+\rho^{-1}\sup_{|x'|<3\rho}u\bigr).
\]

Combining the above estimates, we therefore have

\[
\int_{\mathcal{G}_\rho} w\,\dA \le \int_{|x'|,\,|u|<\rho} wv\,\dxp
\]

\begin{equation}\tag{16.54}
\le C(n)\rho^n(1+H_0R)(1+H_1R^2)^{1/2}\bigl(1+\rho^{-1}\sup_{\Omega}u\bigr).
\end{equation}

Our desired interior gradient bound now follows by combining the estimates
(16.52) (16.53), (16.54) and exponentiating.

\medskip

\noindent\textbf{Theorem 16.5.} \textit{Let $\Omega$ be a domain in $\mathbb{R}^n$ and $u$ a function in $C^2(\Omega)$. Then, for any point $y' \in \Omega$, we have the estimate}

\begin{equation}\tag{16.55}
|Du(y')|\le C_1 \exp\{C_2 \sup_{\Omega}(u-u(y'))/d\},
\end{equation}

\textit{where $d=\dist(y',\partial\Omega)$ and where $C_1=C_1(n,dH_0,d^2H_1)$, $C_2=C_2(n,dH_0,d^2H_1)$.
($H_0$ and $H_1$ being given by (16.51)).}

\medskip

As an immediate consequence of Theorem 16.5, we have the following interior
estimate for non-negative functions.

\medskip

\noindent\textbf{Corollary 16.6.} \textit{Let $\Omega$ be a domain in $\mathbb{R}^n$ and $u$ a non-negative function in $C^2(\Omega)$.
Then for any point $y \in \Omega$, we have the estimate}

\begin{equation}\tag{16.56}
|Du(y')|\le C_1 \exp\{C_2u(y')/d\}
\end{equation}

\textit{where $C_1$, $C_2$ and $d$ are as in Theorem 16.5.}

\medskip

The exponential form of the estimates (16.55) and (16.56) is interesting in that it
cannot be improved. This is evidenced in the case of two dimensional minimal
surfaces by an example in [FN 4]. In the following section, we shall apply Theorem
16.5 to the Dirichlet problem, with continuous boundary values, for the minimal
surface and prescribed mean curvature equations. A further application to the
minimal surface equation is treated by Problem 16.4. Let us conclude this section
by noting the interior estimates for higher order derivatives which now follow
from Theorem 16.5, the Hölder estimate Theorem 12.1 and the Schauder interior
estimates Theorem 6.2 and Problem 6.1.


% PDF page 419
\source{gilbarg-trudinger:page-0419}

\begin{quote}
\textbf{Corollary 16.7.} \emph{Let $\Omega$ be a domain in $\mathbb{R}^n$ and $u$ a function in $C^2(\Omega)$ whose graph has mean curvature $H \in C^k(\Omega)$, $k \geq 1$. Then $u \in C^{k+1}(\Omega)$ and for any point $y \in \Omega$, and multi-index $\beta$, $|\beta| = k + 1$,}
\end{quote}

\begin{equation}
(16.57)\qquad |D^\beta u(y)| \leq C
\end{equation}

\noindent \emph{where $C = C(n, k, |H|_{k,\Omega}, d, \sup |u|)$, $d = \operatorname{dist}(y,\partial\Omega)$.}

\bigskip

\section*{16.3. Application to the Dirichlet Problem}

In this section we study the solvability of the Dirichlet problem with continuous boundary data for both the minimal surface and prescribed mean curvature equations. For the minimal surface equation we have the following extension of Theorem 14.14.

\begin{quote}
\textbf{Theorem 16.8.} \emph{Let $\Omega$ be a bounded $C^2$ domain in $\mathbb{R}^n$. Then the Dirichlet problem $\mathcal{M}u=0$ in $\Omega$, $u=\varphi$ on $\partial\Omega$, is solvable for arbitrary $\varphi \in C^0(\partial\Omega)$ if and only if the mean curvature of the boundary $\partial\Omega$ is everywhere non-negative.}
\end{quote}

\noindent \textbf{Proof.} Let us assume initially that $\partial\Omega \in C^{2,\alpha}$ for some $\alpha > 0$ and that the mean curvature of $\partial\Omega$ is everywhere non-negative. Let $\{\varphi_m\}$ be a sequence of functions in $C^{2,\alpha}(\overline{\Omega})$ which converges uniformly on $\partial\Omega$ to $\varphi$. By Theorem 14.14, the Dirichlet problems $\mathcal{M}u_m=0$ in $\Omega$, $u_m=\varphi_m$ on $\partial\Omega$, are uniquely solvable in $C^{2,\alpha}(\overline{\Omega})$ and from the comparison principle, Theorem 10.1, we have

\begin{equation*}
\sup_{\Omega} |u_{m_1} - u_{m_2}| \leq \sup_{\partial\Omega} |\varphi_{m_1} - \varphi_{m_2}| \to 0 \qquad \text{as } m_1, m_2 \to \infty.
\end{equation*}

Consequently the sequence $\{u_m\}$ converges uniformly on $\Omega$ to some function $u \in C^0(\overline{\Omega})$ with $u = \varphi$ on $\partial\Omega$. Applying Corollary 16.7, together with Arzela's theorem we then obtain $u \in C^2(\Omega)$ and $\mathcal{M}u=0$ in $\Omega$. The result for $C^2$ domains $\Omega$ follows by approximation of $\Omega$ by $C^{2,\alpha}$ domains. The non-existence part of Theorem 16.8 is an immediate consequence of Theorem 14.14. \qed

Existence results for the inhomogeneous equation (16.1) depend on the establishment of an appropriate maximum principle for solutions. By combining Theorem 13.8, Corollary 14.8 and Theorem 15.2, we first have the following basic result for smooth boundary data (see also Theorem 15.15).

\begin{quote}
\textbf{Theorem 16.9.} \emph{Let $\Omega$ be a bounded domain in $\mathbb{R}^n$ with $\partial\Omega \in C^{2,\alpha}$ for some $\alpha$, $0<\alpha<1$, and let $\varphi \in C^{2,\alpha}(\overline{\Omega})$. Let $H \in C^1(\overline{\Omega})$ and suppose that the mean curvature of $\partial\Omega$, $H'$, satisfies}
\end{quote}

\begin{equation}
(16.58)\qquad H'(y) \geq \frac{n}{n-1} |H(y)|
\end{equation}

% PDF page 420
\source{gilbarg-trudinger:page-0420}

\pagestyle{empty}



\setcounter{page}{408}

\noindent\textbf{16. Equations of Mean Curvature Type}\hfill 408

\bigskip

at each point $y \in \partial\Omega$. Then there exists a unique function $u \in C^{2,\alpha}(\overline{\Omega})$ satisfying
equation (16.1) in $\Omega$ and $u=\varphi$ on $\partial\Omega$ provided the family of solutions of the Dirichlet
problems

\begin{equation*}\tag{16.59}
\M u = \sigma nH(x)\bigl(1+|Du|^2\bigr)^{3/2}\ \text{in }\Omega,\qquad u=\sigma\varphi\ \text{on }\partial\Omega,
\end{equation*}

is uniformly bounded in $\Omega$.

\medskip

The necessity of the condition (16.58) is demonstrated by Corollary 14.13. Let us
now determine a further necessary condition for solvability. Namely, by the integral
form (16.42) of equation (16.1), we have for any $\eta \in C_0^1(\Omega)$

\[
\left|\int_\Omega H\eta\,dx\right|\leq \frac{1}{n}\int_\Omega |\nu\cdot D\eta|\,dx
\]
\[
\leq \frac{1}{n}\sup_\Omega \frac{|Du|}{\sqrt{1+|Du|^2}}\int_\Omega |D\eta|\,dx
\]

\noindent and hence, writing

\[
1-\varepsilon_0=\sup_\Omega \frac{|Du|}{\sqrt{1+|Du|^2}},
\]

\noindent we obtain

\begin{equation*}\tag{16.60}
\left|\int_\Omega H\eta\,dx\right|\leq \frac{1-\varepsilon_0}{n}\int_\Omega |D\eta|\,dx,
\end{equation*}

\noindent for all $\eta \in C_0^1(\Omega)$, and some $\varepsilon_0>0$. However, it is clear from the proof of Theorem
10.10, that condition (16.60) is also sufficient to ensure that an a priori estimate for
 $\sup |u|$ holds for $C^1(\overline{\Omega})$ solutions $u$ of equation (16.1). Accordingly we have, from
Theorem 16.9, the following sharp existence theorem.

\medskip

\noindent\textbf{Theorem 16.10.} \textit{Let $\Omega$ be a bounded domain in $\mathbb{R}^n$ with $\partial\Omega \in C^{2,\alpha}$ for some $\alpha$,
$0<\alpha<1$, and let $\varphi \in C^{2,\alpha}(\overline{\Omega})$. Let $H \in C^1(\overline{\Omega})$ satisfy (16.58) and (16.60). Then
the Dirichlet problem $\M u = nH(1+|Du|^2)^{3/2},\ u=\varphi\ \text{on }\partial\Omega,$ is uniquely solvable
for $u \in C^{2,\alpha}(\overline{\Omega})$. Furthermore, if we only assume that $\varphi \in C^0(\partial\Omega)$, the problem is
uniquely solvable for $u \in C^0(\overline{\Omega}) \cap C^2(\Omega)$.}

\medskip

Note that condition (16.60) is implied by the condition

\begin{equation*}\tag{16.61}
\int_\Omega |H|^n\,dx < \omega_n,
\end{equation*}


% PDF page 421
\source{gilbarg-trudinger:page-0421}

(cf. (10.35)). A more general condition than (16.61) is given in [GT 2]. The second
assertion in Theorem 16.10 follows from the first assertion there in the same way
that Theorem 16.8 follows from Theorem 14.14. Note that, if we only have $u \in$
$C^0(\overline{\Omega}) \cap C^2(\Omega)$, then the function $H$ need only satisfy condition (16.60) with
$\varepsilon_0 = 0$.

    When the function $H$ is constant in $\Omega$ it turns out that the condition (16.60)
in Theorem 16.10 is redundant. To show this we assume that (16.58) holds and let
$\Omega_1$ be the subset of $\Omega$ consisting of points having a unique closest point on $\partial\Omega$.
An examination of the proofs of Lemmas 14.16 and 14.17 then shows that

\[
\Delta d \le -n|H|
\]

in $\Omega_1$, where $d(x)=\operatorname{dist}(x,\partial\Omega)$. Let us now set, for $u \in C^0(\overline{\Omega}) \cap C^2(\Omega)$, $x \in \Omega_1$,

\[
 v(x)=\sup_{\partial\Omega}|u|+\frac{e^{\mu\delta}}{\mu}\bigl(1-e^{-\mu d(x)}\bigr),
\]

where $\delta=\operatorname{diam}\Omega$ and $\mu=1+n|H|$. It then follows that

\[
\Mv v = \bigl[-\mu + (1+|Dv|^2)\,\Delta d\bigr]e^{\mu(\delta-d)}
\]
\[
\le -\bigl[\mu+n|H|(1+|Dv|^2)\bigr]e^{\mu(\delta-d)}
\]
\[
\le -n|H|(1+|Dv|^2)^{3/2},
\]

so that the function $v$ is a supersolution of equation (16.1) in the open set $\Omega_1$. Consequently if the function $u$ satisfies (16.1) in $\Omega$, then, by the comparison principle,
Theorem 10.1, $w=u-v$ assumes a maximum value in $\overline{\Omega}$ either on $\partial\Omega$ or in $\Omega-\Omega_1$.
Now let $y$ be a point in $\Omega-\Omega_1$ and $\gamma$ be a straight line segment from $y$ to $\partial\Omega$,
normal to $\partial\Omega$. If the maximum value of $w$ on $\gamma$ is taken on at $y$ we must have
$Du(y)\neq 0$ and also that the maximum value of $u$ on $\gamma$ occurs at $y$. This shows that $w$
cannot take on a maximum value on $\Omega-\Omega_1$. Therefore we obtain the estimate

\begin{equation}
(16.62) \qquad \sup_{\Omega}|u| \le \sup_{\partial\Omega}|u| + (e^{n\delta}-1)/\mu.
\end{equation}

Combining (16.62) with Theorem 16.9 and Corollary 14.13, we thus have the
following existence theorem for the equation of constant mean curvature.

\begin{center}
\textbf{Theorem 16.11.}  \emph{Let $\Omega$ be a bounded $C^2$ domain in $\mathbb{R}^n$. Then the Dirichlet problem}
\[
\Mv u=nH(1+|Du|^2)^{3/2},\qquad u=\varphi \text{ on } \partial\Omega,
\]
\emph{is soluble for constant $H$ and arbitrary $\varphi \in C^0(\partial\Omega)$ if and only if the mean curvature of $\partial\Omega$, $H'$, satisfies $H'\ge n|H|/(n-1)$ everywhere on $\partial\Omega$.}
\end{center}

% PDF page 422
\source{gilbarg-trudinger:page-0422}

\noindent 16.4. \ Equations in Two Independent Variables

So far this chapter has been concerned with the prescribed mean curvature equation
and, in particular, the minimal surface equation. Now, in the case of two indepen-
dent variables, a somewhat more general class of equations will be considered. We
shall consider equations of the form

\begin{equation}\tag{16.63}
Qu=a^{ij}(x,u,Du)D_{ij}u+b(x,u,Du)=0,
\end{equation}

\noindent where $x=(x_1,x_2)\in\Omega$, $\Omega$ is a domain in $\mathbb{R}^2$ and where $a^{ij}$, $b$, $i$, $j=1,2$, denote
given real-valued functions on $\Omega\times\mathbb{R}\times\mathbb{R}^2$ with

\begin{equation}\tag{16.64}
|\xi|^2-\frac{(p\cdot\xi)^2}{1+|p|^2}\le a^{ij}(x,z,p)\xi_i\xi_j\le \gamma\left[|\xi|^2-\frac{(p\cdot\xi)^2}{1+|p|^2}\right]
\end{equation}

\noindent for all $(x,z,p)\in\Omega\times\mathbb{R}\times\mathbb{R}^2$ and all $\xi=(\xi_1,\xi_2)\in\mathbb{R}^2$; and

\begin{equation}\tag{16.65}
|b(x,z,p)|\le \mu\sqrt{1+|p|^2}
\end{equation}

\noindent for all $(x,z,p)\in\Omega\times\mathbb{R}\times\mathbb{R}^2$. Here $\gamma$ and $\mu$ denote fixed constants.
Note that the minimal surface equation, $\mathfrak{M}u=0$, can be written in the form
(16.63) with

\begin{equation*}
a^{ij}(x,z,p)=\delta^{ij}-\frac{p_ip_j}{1+|p|^2},\qquad b=0;
\end{equation*}

\noindent in this case (16.64) and (16.65) hold with $\gamma=1$ and $\mu=0$. More generally, any
equation which arises as the non-parametric Euler equation of an elliptic parametric
functional (see Appendix 16.8), is of the form (16.63), (16.64), (16.65). But quite
apart from these examples, the class of equations (16.63), (16.64), (16.65), which we
call the class of equations of mean curvature type, is both natural and interesting in
that it is completely characterized as follows:

\noindent\hspace*{2em}Suppose $u$ is a $C^2(\Omega)$ function with graph

\[
\mathcal{G}=\{(x,z)\in\mathbb{R}^3\mid x\in\Omega,\ z=u(x)\}.
\]

\noindent Then there exist real-valued functions $a^{ij}$, $b$ such that (16.63), (16.64), (16.65) hold
if and only if the principal curvatures $\kappa_1,\kappa_2$ of $\mathcal{G}$ (see Section 14.6) are related at each
point of $\mathcal{G}$ by an equation of the form

\begin{equation}\tag{16.63$'$}
\alpha_1\kappa_1+\alpha_2\kappa_2+\beta=0,
\end{equation}

\noindent with $\alpha_1$, $\alpha_2$, $\beta$ satisfying

\begin{equation}\tag{16.64$'$}
1\le \alpha_i\le \gamma,\qquad i=1,2,
\end{equation}

\begin{equation}\tag{16.65$'$}
|\beta|\le \mu.
\end{equation}


% PDF page 423
\source{gilbarg-trudinger:page-0423}

To demonstrate that this characterization is valid, we let $d$ denote the distance
function of $\Ssurface$ defined for $X=(x,z)\in \Omega\times \R$ by setting
$d(X)=\dist(X,\Ssurface)$ if $z>u(x)$ and
$d(X)=-\dist(X,\Ssurface)$ if $z<u(x)$. Since $d$ is $C^2$ (see Lemma 14.16) and
$d(x,u(x))\equiv 0$, $x\in \Omega$, we then have, by the chain rule, the identities

\[
D_i d(X)+D_i u(x)D_3 d(X)=0
\]

and

\begin{equation}
D_{ij}d(X)+D_i u(x)D_3 d(X)+D_j u(x)D_{3i}d(X)+D_i u(x)D_j u(x)D_{33}d(X)
+D_3 d(X)D_{ij}u(x)=0,
\tag{16.66}
\end{equation}

$i,j=1,2$, where $X=(x,u(x))$. Since $D_3 d(X)=v^{-1}$, $v=\sqrt{1+|Du(x)|^2}$, (16.63)
then implies

\begin{equation}
\sum_{i,j=1}^3 a_*^{ij}(x)D_{ij}d(X)+v^{-1}b_*(x)=0,
\tag{16.67}
\end{equation}

where $b_*(x)=b(x,u(x),Du(x))$ and where the $3\times 3$ matrix $[a_*^{ij}(x)]$ is defined by
setting $a_*^{ij}(x)=a^{ij}(x,u(x),Du(x))$ for $i,j=1,2$, and

\[
a_*^{i3}(x)=a_*^{3i}(x)=\sum_{j=1}^2 D_j u(x)a_*^{ij}(x),\qquad i=1,2,
\]
\[
a_*^{33}(x)=\sum_{i,j=1}^2 D_i u(x)D_j u(x)a_*^{ij}(x).
\]

Note that these last relations are equivalent to

\[
\sum_{j=1}^3 a_*^{ji}(x)v_j=\sum_{j=1}^3 a_*^{ij}(x)v_j=0,\qquad i=1,2,3,
\]

where $\nu=v^{-1}(-Du(x),1)(=Dd(X))$ is the upward unit normal of $\Ssurface$. Next we let q
be the matrix of an isometry which transforms the coordinate system to a principal
coordinate system at $X$ (see Section 14.6), so that
$q^t[D_{ij}d(X)]q=\operatorname{diag}[\kappa_1,\kappa_2,0]$,
where $\kappa_1,\kappa_2$ are principal curvatures of $\Ssurface$ at $X$. Thus (16.67) can be written in the
form (16.63)${}'$, with $\alpha_1,\alpha_2$ the first two elements or the leading diagonal of
$q^t[a_*^{ij}(x)]q$
and with $\beta=v^{-1}b_*(x)$. (16.65)${}'$ is now true by (16.65). To check (16.64)${}'$, we first
note that

\[
\sum_{i,j=1}^3 a_*^{ij}(x)\xi_i\xi_j=\sum_{i,j=1}^2 a_*^{ij}(x)(\xi_i+\xi_3D_i u(x))(\xi_j+\xi_3D_j u(x)),\qquad \xi\in \R^3,
\]


% PDF page 424
\source{gilbarg-trudinger:page-0424}

and it then follows from (16.64) that
\[
|\xi'|^2 \le \sum_{i,j=1}^3 a_*^{ij}(x)\xi_i\xi_j \le \gamma|\xi'|^2, \qquad \xi' = \xi - (\nu\cdot \xi)\nu,\ \nu = v^{-1}(-Du(x),1).
\]
(16.64)\,' now easily follows from this.

To prove the converse implication we suppose that (16.63)\,' (16.64)\,' and (16.65)\,' hold at $X=(x,u(x))\in\G$, we let $[a_*^{ij}(x)] = q\,\mathrm{diag}\,[\alpha_1,\alpha_2,0]q$, where $q$ is as above, and we let $b_*(x)=v\beta$. Then (16.67) holds and consequently, since we still have the relations
\[
\sum_{j=1}^3 a_*^{ij}(x)\nu_j=0, \qquad i=1,2,3,
\]
an application of (16.66) yields
\[
\sum_{i,j=1}^2 a_*^{ij}(x)D_{ij}u+b_*(x)=0.
\]

We then define, for $i,j=1,2$,
\[
a^{ij}(x,z,p)=\begin{cases}
 a_*^{ij}(x) & \text{if } z=u(x)\text{ and }p=Du(x)\\
 \delta_{ij}-\dfrac{p_ip_j}{1+|p|^2} & \text{otherwise,}
\end{cases}
\]
and
\[
b(x,z,p)=\begin{cases}
 b_*(x) & \text{if } z=u(x)\text{ and }p=Du(x)\\
 0 & \text{otherwise.}
\end{cases}
\]

(16.63), (16.64), (16.65) are now easily checked.
The treatment of equations of the form (16.63), (16.64), (16.65), to be given here,
is in many respects analogous to the treatment given in Chapter 12 for uniformly
elliptic equations in two independent variables; as in Chapter 12 we shall begin by
considering quasiconformal maps, although here it will be necessary to consider
mappings between surfaces in $\R^3$ rather than mappings in the plane. Also, as in
Chapter 12, the principal result is a H"older estimate for quasiconformal maps. A
special consequence of this general estimate is a H"older estimate for the unit normal
of the graph of a solution $u$ of (16.63), (16.64), (16.65). Using this estimate, apriori
bounds for the principal curvatures of the graph of $u$ and for the gradient of $u$ will
be established. One of the most striking results of the theory of equations of the
form (16.63), (16.64), (16.65) is the following generalization of a classical theorem
of Bernstein:

\textbf{Theorem 16.12.} \textit{Let $u\in C^2(\R^2)$ satisfy equation (16.63) on the whole of $\R^2$, with
$b\equiv 0$, and suppose that (16.64) holds with $\Omega=\R^2$. Then $u$ is a linear function.}


% PDF page 425
\source{gilbarg-trudinger:page-0425}

\section*{16.5. Quasiconformal Mappings}

We shall see below that this theorem is also a consequence of the H\"older estimate for the unit normal of the graph of $u$.

\section*{16.5. Quasiconformal Mappings}

In this section, we shall consider mappings between $C^2$ hypersurfaces $\Sc,\Tc\subset\mathbb{R}^3$.
The surfaces $\Sc$ and $\Tc$ are assumed to be \emph{oriented}, that is, it is assumed that there
exist unit normal vectors $\nu,\mu$ defined and continuous over all of $\Sc,\Tc$ respectively.
For our application in the next section $\Sc$ and $\Tc$ will be graphs and hence given by
representations of the form (16.2), (16.3).

Points of $\mathbb{R}^3$ will be denoted $X=(x_1,x_2,x_3)$, $Y=(y_1,y_2,y_3)$ will always denote a
fixed point of $\Sc$ and, for $\rho>0$, we set $\Sc_\rho(Y)=\Sc\cap B_\rho(Y)$. $R$ will denote a fixed
positive constant such that $\Sc_R(Y)\subset\subset\Sc$. We shall often use $\Sc_\rho$ as an abbreviation
for $\Sc_\rho(Y)$.

We shall need the classical version of Stokes' theorem: if $v=(v_1,v_2,v_3)$ is a $C^1$
vector function defined in a neighborhood of $\Sc$ and if $\Gc\subset\subset\Sc$ is such that $\partial\Gc$
consists of a finite union of simple closed $C^1$ curves, then

\begin{equation}\tag{16.68}
\int_{\Gc} v\cdot \mathrm{curl}\, v\, dA \equiv \int_{\Gc} v\cdot D\times v\, dA
= \int_{\partial\Gc} v_i\,dx_i \equiv \int_{\partial\Gc} \mathbf{t}\cdot v\, ds,
\end{equation}

where $A$ denotes surface area on $\Sc$, $\partial\Gc$ is appropriately oriented, $s$ denotes arc
length on $\partial\Gc$ and $\mathbf{t}$ is the unit tangent vector of $\partial\Gc$. In this section we follow the
convention that repeated indices imply summation from 1 to 3. If we take $v=fDg$
in (16.68), where $f,g$ are respectively $C^1$ and $C^2$ functions defined in a neighborhood
of $\Sc$, then, by virtue of the operator identity $D\times D=0$, we get from (16.68)

\begin{equation}
\int_{\Gc} v\cdot Df\times Dg\, dA = \int_{\partial\Gc} f\,dg \equiv \int_{\partial\Gc} f\,\frac{dg}{ds}\, ds,
\end{equation}

where $dg/ds\equiv \mathbf{t}\cdot Dg$ is the directional derivative of $g$ in the direction of $\mathbf{t}$. Since only
first derivatives of $f,g$ appear here, it is easy to see that the above identity is valid
if both $f$ and $g$ are merely $C^1$. Note that, because of the vector identities $a\times a=0$,
$a\cdot a\times b=0$ for $a,b\in\mathbb{R}^3$, we can write

\[
v\cdot Df\times Dg = v\cdot \delta f\times \delta g
\]

on $\Sc$, where $\delta$ is the tangential gradient operator on $\Sc$ defined by (16.4). Hence
for arbitrary $C^1$ functions $f,g$ on $\Sc$ we have the identity

\begin{equation}\tag{16.69}
\int_{\Gc} v\cdot \delta f\times \delta g\, dA = \int_{\partial\Gc} f\,dg \equiv \int_{\partial\Gc} f\,\frac{dg}{ds}\, ds.
\end{equation}


% PDF page 426
\source{gilbarg-trudinger:page-0426}

Our basic assumption concerning $\Tc$ will be that there exists a vector function
$\omega=(\omega_1,\omega_2,\omega_3)$ which is $C^1$ in some neighborhood of $\Tc$ and such that

\begin{equation}\tag{16.70}
\sup_{\Tc}|\omega|+\sup_{\Tc}|D\omega|\leq \Lambda_0\qquad \text{and}\qquad \mu\cdot D\times \omega \equiv 1\text{ on }\Tc.
\end{equation}

where $\Lambda_0$ is a constant. By applying Stokes' formula (16.68) on a subset $\Gc\subset\Tc$
(with $\mu$ in place of $v$ and $\omega$ in place of $v$), we then have

\begin{equation}\tag{16.71}
A(\Gc)=\int_{\partial\Gc}\omega_i\,dx_i,
\end{equation}

that is, the area of a subset $\Gc\subset\Tc$ can be expressed as a boundary integral taken
over $\partial\Gc$.

An example of special interest for us here will be the case when $\Tc$ is the upper
hemisphere $\{X=(x_1,x_2,x_3)\mid |X|=1,\ x_3>0\}$ of the unit sphere. In this case,
taking $\mu(X)=X$, we have (16.70) with

\[
\omega(X)=(-x_2/(1+x_3),\ x_1/(1+x_3),\ 0).
\]

It is worth pointing out (although we shall have no need of it here) that if $\Omega$ is an
arbitrary, connected, oriented, compact $C^2$ surface in $\mathbb{R}^3$, if $\Rc\subset\Omega$ is compact
with non-empty interior and if we take

\[
\Tc=\Omega-\Rc,
\]

then there always exists a vector field $\omega$ satisfying (16.70). A proof of this assertion
involves a straightforward application of the theory of differential forms and the
de Rham cohomology groups; the reader is referred to the discussion in [SI 6].
Let us now consider a mapping

\[
\varphi=(\varphi_1,\varphi_2,\varphi_3):\Gc\to\Tc
\]

which is $C^1$ in the sense that each $\varphi_i$ (as a mapping from $\Gc$ into $\mathbb{R}$) has a $C^1$ extension
$\tilde\varphi_i$ to some neighborhood of $\Gc$. We wish to introduce the concept of quasiconfor-
mality of $\varphi$, but to do this we first need to define the signed area magnification factor
$J(\varphi)$ of $\varphi$. Namely, $J(\varphi)$ is defined on $\Gc$ by

\begin{equation}\tag{16.72}
J(\varphi)=v\cdot \delta(\omega_i\circ \varphi)\times (\delta\varphi_i).
\end{equation}

This definition of $J(\varphi)$ is easily motivated as follows. Let $\Ec$ be a region of $\Gc$ which is
such that $\partial\Ec$ is a simple, smooth curve, and suppose that $\varphi$ is one-to-one with a $C^1$
inverse in some open subset of $\Tc$ which contains $\varphi(\Ec)$. Assuming that the curves
$\partial\Ec$ and $\partial\varphi(\Ec)$ are appropriately oriented, we can apply (16.69) and (16.71) to give

\[
A(\varphi(\Ec))=\int_{\partial\varphi(\Ec)}\omega_i\,dx_i=\pm\int_{\partial\Ec}\omega_i\circ\varphi\,d\varphi_i=\pm\int_{\Ec} v\cdot (\delta\omega_i\circ\varphi)\times(\delta\varphi_i)\,dA
\]


% PDF page 427
\source{gilbarg-trudinger:page-0427}

with the $+$ or $-$ sign according as $\varphi$ is orientation preserving or reversing on $\Sc$.
This identity clearly motivates the definition (16.72).

An important (and intuitively obvious) fact concerning $J(\varphi)$ is that it is independent of coordinates in the following sense: If $P, Q$ are linear isometries of $\mathbb{R}^3$ with

\[
\det P=\det Q=1,
\]

and if we define

\[
\widetilde{\Sc}=\{P(X-Y)\mid X\in\Sc\},\ \widetilde{\Tc}=\{Q(X-\varphi(Y))\mid X\in\Tc\},
\]
\[
\widetilde{v}=P\circ v\circ P^{-1},\ \widetilde{\mu}=Q\circ \mu\circ Q^{-1},\ \widetilde{\omega}=Q\circ \omega\circ Q^{-1},\ \widetilde{\varphi}=Q\circ \varphi\circ P^{-1},
\]

then

\begin{equation}\tag{16.73}
\mathbf{v}\cdot(\widetilde{\delta}\omega_i\circ\varphi)\times(\delta\varphi_i)(Y)=\widetilde{\mathbf{v}}\cdot(\widetilde{\delta}\widetilde{\omega}_i\circ\widetilde{\varphi})\times(\widetilde{\delta}\widetilde{\varphi}_i)(0),
\end{equation}

where $\widetilde{\delta}$ denotes the tangential gradient operator on $\widetilde{\Sc}$. The relation (16.73) is
easily checked by first representing the isometries $P, Q$ in terms of orthogonal
matrices and then using two elementary facts from linear algebra, namely that if
$A, B$ are $3\times 3$ matrices, then $\det AB=\det A\,\det B$, and if $A$ has rows $a, b, c$, then
$\det A=a\cdot b\times c$. It can also be checked by using (16.70) that

\[
\widetilde{\mu}\cdot D\times \widetilde{\omega}\equiv 1\quad \text{on }\widetilde{\Sc}.
\]

If the isometries $P, Q$ above are chosen so that

\[
\widetilde{v}(0)=\widetilde{\mu}(0)=(0,0,1),
\]

and if we introduce new coordinates

\[
(\zeta,\zeta_3)=(\zeta_1,\zeta_2,\zeta_3)=P(X-Y),
\]

then $\widetilde{\Tc}$ can be represented near $0$ in the form

\[
\zeta_3=\widetilde{u}(\zeta),\ \zeta\in\Uc,
\]

where $\Uc$ is a neighborhood of $0\in\mathbb{R}^2$ and $\widetilde{u}$ is a $C^2(\Uc)$ function with
$D\widetilde{u}(0)=0$.
We also then have

\[
\widetilde{\delta}_j\widetilde{\varphi}_3(0)=(0,0,1)\cdot\widetilde{\delta}_j\widetilde{\varphi}(0)
=\widetilde{\mu}(0)\cdot\widetilde{\delta}_j\widetilde{\varphi}(0)=0,\qquad j=1,2,3
\]

by virtue of the fact that the vectors $\widetilde{\delta}_j\widetilde{\varphi}(0)$, $j=1,2,3$, are tangent to $\widetilde{\Tc}$ at $0$. That is,
we have

\[
\widetilde{\delta}\widetilde{\varphi}_3(0)=0.
\]


% PDF page 428
\source{gilbarg-trudinger:page-0428}

Thus if
\[
\psi : \mathcal U \to \mathbb R^2
\]
is defined by
\[
\psi(\zeta) = [\tilde{\varphi}_1(\zeta,\tilde{u}(\zeta)), \tilde{\varphi}_2(\zeta,\tilde{u}(\zeta))], \qquad \zeta \in \mathcal U,
\]
then $\psi$ approximates $\tilde{\varphi}$ near $0$ in the sense that
\[
\bigl|(\psi(\zeta),0) - \tilde{\varphi}(\zeta,\tilde{u}(\zeta))\bigr| = o(|\zeta|) \text{ as } |\zeta| \to 0, \qquad \zeta \in \mathcal U.
\]

Furthermore, using (16.73) together with the definition $\tilde{\delta} = D - \tilde{v}(\tilde{v} \cdot D)$, one now easily checks that

\begin{equation}
J(\varphi)(Y) = D_1\psi_1(0)D_2\psi_2(0) - D_1\psi_2(0)D_2\psi_1(0);
\tag{16.74}
\end{equation}

that is, $J(\varphi)(Y)$ is just the Jacobian of $\psi$ at $0$. Also it follows that

\begin{equation}
|\delta\varphi(Y)|^2 = |\tilde{\delta}\tilde{\varphi}(0)|^2 = |D\psi(0)|^2.
\tag{16.75}
\end{equation}

In view of (16.74), (16.75) it is now reasonable, by analogy with (12.2), to make the following definition. Namely, the mapping $\varphi$ is said to be a $(K, K')$-quasiconformal mapping from $\mathcal G$ into $\mathcal Z$ if

\begin{equation}
|\delta\varphi(X)|^2 \leq 2KJ(\varphi)(X) + K'.
\tag{16.76}
\end{equation}

at each point $X \in \mathcal G$. Here $K, K'$ are real constants with $K' \geq 0$ and $|\delta\varphi(X)|^2 = \sum_{i=1}^3 |\delta\varphi_i(X)|^2$. We emphasize here that it is not assumed in the above that $K$ is positive (cf. (12.2)). Note that when $K' = 0$, we must have $|K| \geq 1$ unless $\delta\varphi \equiv 0$.

Before beginning the proof of the main Hölder continuity result for quasiconformal maps, one further preliminary is needed; namely, if $g$ is an arbitrary $C^1$ function on $\mathcal G_\rho$, then

\begin{equation}
\int_{\partial \mathcal G_\rho} g\, ds = D_\rho \int_{\mathcal G_\rho} g|\delta r|\, dA
\tag{16.77}
\end{equation}

for almost all $\rho \in (0,R)$, where $r$ is the radial distance function relative to $Y$, defined by
\[
r(X) = |X - Y|, \qquad X \in \mathbb R^3.
\]

Formula (16.77) is in fact a special case of the important co-area formula (see [FE]). We proceed to prove (16.77). First note that the left side of (16.77) makes sense for

% PDF page 429
\source{gilbarg-trudinger:page-0429}

\pagestyle{empty}


almost all $\rho \in (0, R)$ because by Sard's Theorem [SB] we can write

\begin{equation}\tag{16.78}
\partial \widetilde{\mathscr{S}}_{\rho}
= \widetilde{\mathscr{S}}_{\rho} \cap \{X \in \mathbb{R}^{3}\mid |X-Y|=\rho\}
= \bigcup_{j=1}^{n(\rho)} \Gamma_{\rho}^{(j)}
\end{equation}

for almost all $\rho \in (0, R)$, where the $\Gamma_{\rho}^{(j)}$ are simple closed $C^{2}$ curves and $n(\rho)$ is a positive integer. Actually, Sard's Theorem guarantees that for almost all $\rho \in (0, R)$ the tangential gradient $\delta r$ vanishes at no point of $\partial \widetilde{\mathscr{S}}_{\rho}$; the geometric interpretation of this is that the surface $\widetilde{\mathscr{S}}$ and the sphere $\{X \in \mathbb{R}^{3}\mid |X-Y|=\rho\}$ intersect non-tangentially, and this explains why (16.78) holds. Now let us take some fixed $\rho \in (0, R)$ with $\delta r \neq 0$ on $\partial \widetilde{\mathscr{S}}_{\rho}$ and let $\varepsilon > 0$ be small enough to ensure $\delta r \neq 0$ on $\partial \widetilde{\mathscr{G}}$, where $\widetilde{\mathscr{G}} = \widetilde{\mathscr{S}}_{\rho} - \widetilde{\mathscr{S}}_{\rho-\varepsilon}$. We wish now to apply Stokes' Theorem (16.68); the appropriately oriented unit tangent for $\partial \widetilde{\mathscr{G}}$ is given by $\mathbf{t}=\mathbf{v}\times \delta r/|\delta r|$ on $\partial \widetilde{\mathscr{S}}_{\rho}$ and $\mathbf{t}=-\mathbf{v}\times \delta r/|\delta r|$ on $\partial \widetilde{\mathscr{S}}_{\rho-\varepsilon}$. Let $\mathbf{F}$ denote a $C^{1}$ extension of the vector function $\mathbf{v}\times \delta r/|\delta r|$ to some open subset of $\mathbb{R}^{3}$ containing $\widetilde{\mathscr{G}}$. Then applying (16.68) with

\[
\mathbf{v}=\frac{1}{\varepsilon}(r-\rho+\varepsilon)g\mathbf{F},
\]

and noting that $\mathbf{v}=0$ on $\partial \widetilde{\mathscr{S}}_{\rho-\varepsilon}$, we obtain

\[
\int_{\partial \widetilde{\mathscr{S}}_{\rho}} g\,ds
= \int_{\partial \widetilde{\mathscr{G}}} \frac{1}{\varepsilon}(r-\rho+\varepsilon)g\mathbf{F}\cdot \mathbf{t}\,ds
\]

\[
= \int_{\widetilde{\mathscr{S}}_{\rho}-\widetilde{\mathscr{S}}_{\rho-\varepsilon}}
\frac{1}{\varepsilon}(r-\rho+\varepsilon)\mathbf{v}\cdot(gD\times \mathbf{F}+Dg\times \mathbf{F})\,dA
\]

\[
+ \frac{1}{\varepsilon}\int_{\widetilde{\mathscr{S}}_{\rho}-\widetilde{\mathscr{S}}_{\rho-\varepsilon}}
g\,\mathbf{v}\cdot D r\times \mathbf{F}\,dA.
\]

Since

\[
\mathbf{v}\cdot D r\times \mathbf{F}=(\mathbf{v}\times D r)\cdot(\mathbf{v}\times \delta r/|\delta r|)
=|\mathbf{v}\times \delta r|^{2}/|\delta r|=|\delta r|
\]

on $\widetilde{\mathscr{S}}_{\rho}-\widetilde{\mathscr{S}}_{\rho-\varepsilon}$, we then have

\[
\int_{\partial \widetilde{\mathscr{S}}_{\rho}} g\,ds
= \int_{\widetilde{\mathscr{S}}_{\rho}-\widetilde{\mathscr{S}}_{\rho-\varepsilon}}
\frac{1}{\varepsilon}(r-\rho+\varepsilon)\mathbf{v}\cdot(gD\times \mathbf{F}+Dg\times \mathbf{F})\,dA
\]

\[
+ \frac{1}{\varepsilon}\left\{\int_{\widetilde{\mathscr{S}}_{\rho}} g|\delta r|\,dA
- \int_{\widetilde{\mathscr{S}}_{\rho-\varepsilon}} g|\delta r|\,dA\right\}.
\]


% PDF page 430
\source{gilbarg-trudinger:page-0430}

and, since

\[
0 \le \frac{r-\rho+\varepsilon}{\varepsilon} \le 1
\]

on $\Sscr_{\rho}-\Sscr_{\rho-\varepsilon}$, we have (16.77) by letting $\varepsilon \to 0$.
The main H\"older estimate will be a consequence of estimates for the integral
\[
\D(\rho; Z),
\]
which is defined for $Z \in \Gscr$ and $\rho > 0$ by
\[
\D(\rho; Z)=\int_{\Sscr_{\rho}(Z)} |\delta \varphi|^2\, dA.
\]

(cf. the Dirichlet integral used in Chapter 12).
The following lemma, which should be compared with inequality (12.8),
provides a bound for $\D(R/2; Y)$. In the statement of the lemma, and subsequently,
$\Lambda_1$ denotes a constant such that

\begin{equation}\tag{16.79}
\A(\Sscr_{3R/4}(Y)) \le \Lambda_1 (3R/4)^2.
\end{equation}

It will be shown in the next section that in the case where $\Gscr$ is a graph with $(K,K')$-
quasiconformal Gauss map, one can obtain a bound for $\Lambda_1$ in terms of $K$ and $K'R^2$.

\medskip

\noindent\textbf{Lemma 16.13.}\quad \textit{Suppose $\varphi$ is a $(K,K')$-quasiconformal $C^1$ mapping from $\Gscr$ into $\Tscr$.
Then}

\begin{equation}\tag{16.80}
\D(R/2; Y) \le C
\end{equation}

\noindent\textit{where $C = C(\Lambda_0, K, K'R^2, \Lambda_1)$.}

\medskip

\noindent\textit{Proof.}\quad Let $\Gamma_{\rho}^{(j)}$, $j=1,\ldots,n(\rho)$, be as in (16.78), and assume $\Gamma_{\rho}^{(j)}$ are oriented
appropriately for the use of (16.68) over $\Sscr_{\rho}$. Then, using the definition (16.72) of
$\J(\varphi)$ in combination with (16.69), we obtain the identity

\begin{equation}\tag{16.81}
\int_{\Sscr_{\rho}} \J(\varphi)\, dA
=
\sum_{j=1}^{n(\rho)}
\int_{\Gamma_{\rho}^{(j)}} \omega_i \circ \varphi \,\frac{d\varphi_i}{ds}\, ds
\end{equation}

for almost all $\rho \in (0,R)$. This identity will play a key role in the proof of the main
estimate for $\D(\rho; Z)$ given in the next theorem. For the present, we simply use
Schwarz's inequality in combination with the inequality

\[
\left|\frac{d\varphi_i}{ds}\right| \le |\delta \varphi_i|,\quad i=1,2,3,
\]


% PDF page 431
\source{gilbarg-trudinger:page-0431}

whereupon (16.81), (16.70) and (16.77) imply that for almost all $\rho \in (0,R)$
\[
\left|\int_{\mathfrak{S}_\rho} J(\varphi)\,dA\right|
\leq \left(\sup |\omega|\right)\int_{\partial S_\rho} |\delta\varphi|\,ds
\]
\[
\leq \Lambda_0\left(\int_{\partial S_\rho} |\delta\varphi|^2\,ds\right)^{1/2}
\left(\int_{\partial S_\rho} ds\right)^{1/2}
\]
\[
=\Lambda_0\left(D_\rho\int_{\mathfrak{S}_\rho} |\delta r|\,|\delta\varphi|^2\,dA\right)^{1/2}
\left(D_\rho\int_{\mathfrak{S}_\rho} |\delta r|\,dA\right)^{1/2}.
\]

Thus, since $|\delta r|\leq |Dr|=1$, we obtain from (16.76) that for almost all $\rho \in (0,3R/4)$
\[
\Dop(\rho)\leq 2\Lambda_0|K|(f'(\rho)\Dop'(\rho))^{1/2}+\Lambda_1K'R^2.
\]

where we are using $\Dop(\rho)$ as an abbreviation for $\Dop(\rho;Y)$ and where
\[
f(\rho)=\Aop(\tilde S_\rho(Y)).
\]

Note that $f'(\rho)$ and $\Dop'(\rho)$ exist for almost all $\rho\in(0,R)$ because $f(\rho)$ and $\Dop(\rho)$ are non-decreasing in $\rho$. Squaring each side of the above inequality, we then obtain
\[
\Dop^2(\rho)\leq 8(\Lambda_0 K)^2f'(\rho)\Dop'(\rho)+2(K'\Lambda_1R^2)^2.
\]

Now if we let
\[
g(\rho)=f(\rho)+\rho R
\]

(so that $g'(\rho)\geq R$) and let
\[
\Eop(\rho)=\Dop(\rho)+\rho/R,
\]

then it is quite a straightforward matter to show that the previous inequality implies an inequality of the form
\[
\Eop^2(\rho)\leq Cg'(\rho)\Eop'(\rho)
\]

for almost all $\rho\in(R/4,R/2)$, where $C=C(\Lambda_0,K,K'R^2,\Lambda_1)$. This can be written
\[
(16.82)\qquad -\frac{d}{d\rho}\left(\frac{1}{\Eop(\rho)}\right)\geq \frac{1}{Cg'(\rho)}.
\]


% PDF page 432
\source{gilbarg-trudinger:page-0432}

Now, by using the H\"older inequality and the monotonicity of $g$, we have
\[
\left(\frac{R}{4}\right)^2
=
\left(\int_{R/2}^{3R/4} d\rho\right)^2
\]
\[
\leq
\left(\int_{R/2}^{3R/4} g'(\rho)\,d\rho\right)
\left(\int_{R/2}^{3R/4} \frac{d\rho}{g'(\rho)}\right)
\]
\[
\leq
\bigl(g(3R/4)-g(R/2)\bigr)
\int_{R/2}^{3R/4}\frac{d\rho}{g'(\rho)},
\]
so that
\[
\int_{R/2}^{3R/4}\frac{d\rho}{g'(\rho)}
\geq
\frac{(R/4)^2}{g(3R/4)-g(R/2)}
>
\frac{(R/4)^2}{g(3R/4)}.
\]

Hence integrating (16.82) over $(R/2,3R/4)$ and using this last inequality we obtain
\[
\frac{1}{\Eop(R/2)}-\frac{1}{\Eop(3R/4)}
\geq
\frac{(R/4)^2}{C\,g(3R/4)},
\]
so that
\[
\Eop(R/2)\leq \frac{C}{R^2}\bigl[A(\Sscr_{3R/4}(Y))+R^2\bigr].
\]

The desired result (16.80) now follows by using (16.79). $\Box$

The next theorem contains the main estimate for $\Dop(\rho; Z)$. In the statement of the
theorem, and subsequently, $\Lambda_2$ denotes a constant such that
\[
\tag{16.83}
\int_{\Sscr_{R/2}(Y)} H^2\, dA \leq \Lambda_2.
\]
where $H=(\kappa_1+\kappa_2)/2$ is the mean curvature of $\Sscr$.

\medskip

\noindent\textbf{Theorem 16.14.}\quad \textit{Suppose $\varphi$ is a $(K,K')$-quasiconformal $C^1$ mapping from $\Sscr$ into $\Tscr$.
Then}

\begin{equation}\tag{16.84}
\Dop(\rho; Z)\leq C(\rho/R)^\alpha
\end{equation}


% PDF page 433
\source{gilbarg-trudinger:page-0433}
\pagestyle{empty}

for all $Z \in \mathcal G_{R/4}(Y)$ and all $\rho \in (0,R/4)$, where $C$ and $\alpha$ are positive constants depending only on $\Lambda_0$, $K$, $K'R^2$, $\Lambda_1$, and $\Lambda_2$.

\medskip

Proof. \ In the proof we let $\mathcal G_\rho=\mathcal G_\rho(Z)$ and $\mathcal D(\rho)=\mathcal D(\rho;Z)$, where $Z \in \mathcal D_{R/4}(Y)$, and let $r$ denote the radial distance function defined by $r(X)=|X-Z|$, $X \in \mathbb R^3$. We take $\rho \in (0,R/4)$ such that $\delta r$ is never zero on $\partial\mathcal G_\rho$, so that we can assume (16.77) and (16.78). Since the curves $\Gamma_\rho^{(j)}$ are closed, we have

\[
\int_{\Gamma_\rho^{(j)}} \frac{d\varphi_i}{ds}\, ds = 0; \qquad (\Gamma^{(j)}=\Gamma_\rho^{(j)}),
\]

hence, if we let $X_j$ denote the initial point (corresponding to arc length $s=0$) of $\Gamma_\rho^{(j)}$, then the integral on the right hand side of (16.81) can be written

\[
\int_{\Gamma_\rho^{(j)}} \bigl(\omega_i\circ\varphi-\omega_i\circ\varphi(X_j)\bigr)\frac{d\varphi_i}{ds}\, ds.
\]

Hence (16.81) gives

\[
\left|\int_{\partial\mathcal G_\rho} J(\varphi)\, dA\right|
\leq \sum_{j=1}^{n(\rho)}
\left\{
\sup_{\Gamma_\rho^{(j)}}\bigl|\omega\circ\varphi-\omega\circ\varphi(X_j)\bigr|
\int_{\Gamma_\rho^{(j)}} \left|\frac{d\varphi}{ds}\right|\, ds
\right\}
\]

\[
\leq \sum_{j=1}^{n(\rho)}
\left\{
\int_{\Gamma_\rho^{(j)}} \left|\frac{d\omega\circ\varphi}{ds}\right|\, ds
\int_{\Gamma_\rho^{(j)}} \left|\frac{d\varphi}{ds}\right|\, ds
\right\}
\]

\[
\leq \sum_{j=1}^{n(\rho)}
\left\{
\int_{\Gamma_\rho^{(j)}} |\delta\omega\circ\varphi|\, ds
\int_{\Gamma_\rho^{(j)}} |\delta\varphi|\, ds
\right\}
\]

\[
\leq \sup_Z |D\omega|
\sum_{j=1}^{n(\rho)} \left(\int_{\Gamma_\rho^{(j)}} |\delta\varphi|\, ds\right)^2
\]

\[
\leq \Lambda_0 \left(\sum_{j=1}^{n(\rho)} \int_{\Gamma_\rho^{(j)}} |\delta\varphi|\, ds\right)^2
\qquad \text{by (16.70),}
\]

\[
= \Lambda_0 \left(\int_{\partial\mathcal G_\rho} |\delta\varphi|\, ds\right)^2
\]

\[
\leq \Lambda_0 \left(\int_{\partial\mathcal G_\rho} |\delta\varphi|^2|\delta r|^{-1}\, ds\right)
\left(\int_{\partial\mathcal G_\rho} |\delta r|\, ds\right)
\]

by Schwarz's inequality,


% PDF page 434
\source{gilbarg-trudinger:page-0434}

\[
= \Lambda_0\left(D_\rho \int_{\Sset_\rho} |\delta\varphi|^2\,dA\right)\left(D_\rho \int_{\Sset_\rho} |\delta r|^2\,dA\right)\quad \text{by (16.77),}
\]

\[
\leq C\Lambda_0\rho\,\D'(\rho)
\]

by (16.36), where $C=C(\Lambda_1,\Lambda_2)$. Therefore by (16.76) we have

\[
\D(\rho)\leq C|K|\Lambda_0\rho\,\D'(\rho)+C'K'\rho^2
\]

for almost all $\rho\in(0,R/4)$, where $C'=C'(\Lambda_1,\Lambda_2)$ by virtue of (16.34). Now define

\[
\G(\rho)=\D(\rho)+(\rho/R)^2.
\]

It is then not difficult to see that the above inequality implies an inequality of the form

\[
\G(\rho)\leq C\rho\,\G'(\rho)
\]

where $C=C(\Lambda_0,K,K'R^2,\Lambda_1,\Lambda_2)$. This last inequality can be written

\[
\frac{d}{d\rho}\log \G(\rho)\geq (C\rho)^{-1},
\]

and, since $\G(\rho)$ is increasing in $\rho$, we can integrate to obtain

\[
\log(\G(\rho)/\G(R/4))\leq C^{-1}\log(4\rho/R),\qquad \rho\leq R/4.
\]

Thus

\[
\G(\rho)\leq 4^\alpha \G(R/4)(\rho/R)^\alpha,\qquad \rho\leq R/4,
\]

where $\alpha=C^{-1}$. Since $\Sset_{R/4}(Z)\subset\Sset_{R/2}(Y)$ we must have

\[
\G(R/4)\leq \D(R/2;Y)+1/16
\]

and hence the required estimate follows from (16.80). $\square$

Using the extended Morrey estimate, Lemma 16.4, we can now finally deduce
from Theorem 16.14 the H\"older estimate for $(K,K')$-quasiconformal maps.

\medskip

\noindent\textbf{Theorem 16.15.} \textit{Suppose $\varphi$ is a $(K,K')$-quasiconformal $C^1$ mapping from $\G$ into $\Sigma$. Then}

\begin{equation}\tag{16.85}
\sup_{x\in \Sset_\rho^\ast(Y)} |\varphi(X)-\varphi(Y)|\leq C(\rho/R)^\alpha,\qquad \rho\in(0,R/4),
\end{equation}


% PDF page 435
\source{gilbarg-trudinger:page-0435}




where $C$ and $\alpha$ are positive constants depending only on $\Lambda_0$, $K$, $K'R^2$, $\Lambda_1$, $\Lambda_2$ and $\Sg_\rho^*(Y)$ is the component of $\Sg_\rho(Y)$ which contains $Y$.

\noindent\textit{Proof.} Let $Z$ be an arbitrary point of $\Sg_{R/4}(Y)$. By the H\"older inequality and the estimates (16.84) and (16.34) we have

\[
\int_{\Sg_\rho(Z)} |\delta \varphi_i| \, dA \le (CC')^{1/2}\rho (\rho/R)^{\alpha/2}, \qquad \rho \in (0, R/4),\ i=1,2,3,
\]

\noindent where $C$, $\alpha$ are as in (16.84) and $C' = C'( \Lambda_1, \Lambda_2)$. Hence the hypotheses of Lemma 16.4 are satisfied with $K = (CC')^{1/2}$ and $\beta = \alpha/2$. The theorem thus follows.

\bigskip

\section*{16.6.\ Graphs with Quasiconformal Gauss Map}

In this section $\Sg$ will denote the graph $\{(x, z) \in \R^3 \mid x \in \Omega,\, z = u(x)\}$ of a $C^2(\Omega)$ function $u$, $y$ will denote a fixed point of $\Omega$, and it will be assumed that $\Omega$ contains the disc $B_R(Y) = \{x \in \R^2 \mid |x-y| < R\}$. $Y$ will denote the point $(y, y_3)$ ($y_3 = u(y)$) of $\Sg$ and, as in Section 16.4, $\unitv$ will denote the upward unit normal function defined on $\Omega \times \R$ by

\[
\unitv(x,z) \equiv \unitv(x) = \left(-\frac{\D u(x)}{v}, \frac{1}{v}\right), \qquad v = \sqrt{1+|\D u|^2}, \ x \in \Omega, \ z \in \R.
\]

\noindent The Gauss map $G$ of $\Sg$, which maps $\Sg$ into the upper hemisphere

\[
\Zset = \{X=(x_1,x_2,x_3)\in \R^3 \mid |X|=1, \ x_3>0\},
\]

\noindent is defined by setting

\[
(16.86)\qquad G(x,x_3)=\unitv(x,x_3)
\]

\noindent at each point $(x,x_3)\in \Sg$. The remaining notation and terminology are as in Sections 16.4 and 16.5.

\noindent We first want to explicitly obtain the quantities $|\delta \varphi|^2$ and $J(\varphi)$ in the case when $\varphi=G$. This is quite straightforward if we set up new coordinates as in Section 16.5. In this case the function $\psi$ is defined by

\[
\psi(\zeta) = -\left(1 + |\D \uvec(\zeta)|^2\right)^{-1/2}\D \uvec(\zeta), \qquad \zeta \in \Uset,
\]

\noindent so that we obtain, from (16.74) and (16.75),

\[
J(G)(Y)=D_{11}\uvec(0)D_{22}\uvec(0)-(D_{12}\uvec(0))^2
\]


% PDF page 436
\source{gilbarg-trudinger:page-0436}

\begin{quote}
and

\[
|\partial G|^{2}(Y)=(D_{11}\tilde{u}(0))^{2}+2(D_{12}\tilde{u}(0))^{2}+(D_{22}\tilde{u}(0))^{2}.
\]

But then by Section 14.6 we have

\begin{equation}
\begin{aligned}
J(G)(Y)&=\kappa_{1}\kappa_{2},\\
|\partial G|^{2}(Y)&=\kappa_{1}^{2}+\kappa_{2}^{2}.
\end{aligned}
\tag{16.87}
\end{equation}

where $\kappa_{1}, \kappa_{2}$ are the principal curvatures of $\widetilde{\mathcal S}$ at $Y$. The product $\kappa_{1}\kappa_{2}$ appearing in (16.87) is called the Gauss curvature of $\widetilde{\mathcal S}$; it is an extremely important geometric invariant in the study of surfaces.

By using the identities (16.87) and recalling the definition (16.76), we now see that $G$ is $(K,K')$-quasiconformal if and only if at each point of $\widetilde{\mathcal S}$ the principal curvatures $\kappa_{1}, \kappa_{2}$ satisfy

\begin{equation}
\kappa_{1}^{2}+\kappa_{2}^{2}\leq 2K\kappa_{1}\kappa_{2}+K'.
\tag{16.88}
\end{equation}

We thus see why the study of quasiconformal maps is relevant to the investigation of equations of mean curvature type; because by squaring (16.63)' we obtain the inequality

\[
\frac{\alpha_{1}}{\alpha_{2}}\kappa_{1}^{2}+\frac{\alpha_{2}}{\alpha_{1}}\kappa_{2}^{2}=-2\kappa_{1}\kappa_{2}+\frac{\beta^{2}}{\alpha_{1}\alpha_{2}}.
\]

That is, by virtue of (16.64)' and (16.65)', we deduce that the Gauss map $G$ of the graph of a solution $u$ of (16.63), (16.64), (16.65), is $(K,K')$-quasiconformal with

\begin{equation}
K=-\gamma,\qquad K'=\gamma\mu^{2}.
\tag{16.89}
\end{equation}

Thus the results established in this section for graphs with quasiconformal Gauss map are all applicable to the graphs of solutions of (16.63), (16.64), (16.65).

We wish to eventually apply Theorem 16.15 to the Gauss map $G$, but first we need to discuss appropriate choices for the constants $\Lambda_{0}, \Lambda_{1}$ and $\Lambda_{2}$. To begin with, we have already seen in Section 16.5 that we can take

\[
\omega(X)=\left(-\frac{x_{2}}{1+x_{3}},\frac{x_{1}}{1+x_{3}},0\right).
\]

One then easily checks that an appropriate choice for the constant $\Lambda_{0}$ is $\Lambda_{0}=4$.
Next we notice by Lemma 16.13 and (16.87) we have, provided $G$ is $(K,K')$-quasiconformal,

\[
\int_{\mathfrak{S}_{R/2}(Y)}(\kappa_{1}^{2}+\kappa_{2}^{2})\,dA\leq C,
\]
\end{quote}

% PDF page 437
\source{gilbarg-trudinger:page-0437}




where $C = C(K, K'R^2, \Lambda_1)$. Thus, since

\[
\kappa_1^2+\kappa_2^2 \geq \frac12(\kappa_1+\kappa_2)^2 = 2H^2,
\]

we can take $\Lambda_2 = C/2$, with $C$ as above. The next lemma shows that we can choose
\(\Lambda_1\) to depend only on $K$ and $K'R^2$.

\medskip

\noindent\textbf{Lemma 16.16.} \textit{Suppose $G$ is $(K, K')$-quasiconformal. Then}

\begin{equation}
\tag{16.90}
|\Xset_{\rho/2}(Z)| \leq C\rho^2
\end{equation}

\textit{for each $Z \in \Sset$ and $\rho > 0$ such that $\Xset_{\rho}(Z) \subset \Xset_{R}(Y)$, where $C = C(K, K'R^2)$.}

\medskip

\noindent\textit{Proof.} The starting point for our proof is the identity (16.44) which, by virtue of formula (14.104), can be written in the form

\[
\int_{\Omega} (\kappa_1^2+\kappa_2^2)\eta\,dx + \int_{\Omega} (v_k D_k v_i D_i\eta - 2v_k D_k H\eta)\,dx = 0
\]

for all $\eta \in C_0^1(\Omega)$. Here we write $H(x) = H(x, u(x))$, $\kappa_i(x)=\kappa_i(x, u(x))$ for $x \in \Omega$.
If $\eta \in C_0^2(\Omega)$, we then obtain by integration by parts

\[
\int_{\Omega} (\kappa_1^2+\kappa_2^2-4H^2)\eta\,dx = \int_{\Omega} (v_i v_k D_{ik}\eta - 4Hv_i D_i\eta)\,dx.
\]

and hence, as $H=(\kappa_1+\kappa_2)/2$,

\[
-2\int_{\Omega}\kappa_1\kappa_2\eta\,dx = \int_{\Omega} (v_i v_k D_{ik}\eta - 2(\kappa_1+\kappa_2)v_i D_i\eta)\,dx
\]

for all $\eta \in C_0^2(\Omega)$. By (16.88) we then have, replacing $\eta$ by $\eta^2$,

\[
\int_{\Omega} (\kappa_1^2+\kappa_2^2)\eta^2\,dx \leq |K|\int_{\Omega} \bigl(|D^2\eta^2|+2|\kappa_1+\kappa_2|\,|D\eta^2|\bigr)\,dx + K'\int_{\Omega}\eta^2\,dx.
\]

Since we can write

\[
4K(\kappa_1+\kappa_2)\eta |D\eta|
\leq \frac12(\kappa_1^2+\kappa_2^2)\eta^2 + (4KD\eta)^2,
\]

this gives

\[
\frac12\int_{\Omega}(\kappa_1^2+\kappa_2^2)\eta^2\,dx \leq \int_{\Omega} \{C(|D\eta|^2+\eta|D^2\eta|)+K'\eta^2\}\,dx.
\]


% PDF page 438
\source{gilbarg-trudinger:page-0438}




where $C = C(K)$. Now let $\rho > 0$ and $Z = (z, u(z))$ be such that $B_{\rho}(z) \subset \Omega$, and let us choose $\eta$ such that $0 \leq \eta \leq 1$ in $\Omega$, $\eta = 1$ on $B_{\rho/2}(z)$, $\eta = 0$ on $\Omega - B_{\rho}(z)$, $|D\eta| \leq c/\rho$, $|D^2\eta| \leq c/\rho^2$ where $c$ is an absolute constant. (It is clear that such a function $\eta$ exists.) Then, since $K' \leq K'R^2/\rho^2$, we obtain

\begin{equation}
\tag{16.91}
\int_{\Omega} (\kappa_1^2 + \kappa_2^2)\eta^2\,dx \leq C
\end{equation}

where $C = C(K, K'R^2)$. Consequently, using H\"older's inequality, we have

\[
\int_{\Omega} |H\eta|\,dx \leq \left(\int_{\Omega} (H\eta)^2\,dx\right)^{1/2} |B_{\rho}(z)|^{1/2}
\leq (C\pi)^{1/2}\rho.
\]

The lemma now follows from (16.53). $\square$

From Lemma 16.16, Theorem 16.15 and our previous choices of $\Lambda_0$, $\Lambda_1$, we can now deduce that if $G$ is $(K, K')$-quasiconformal then

\begin{equation}
\tag{16.92}
\sup_{x \in \mathfrak{S}_{\rho}^*(Y)} |\nu(X) - \nu(Y)| \leq C(\rho/R)^{\alpha}, \qquad \rho \in (0, R),
\end{equation}

where $C$ and $\alpha$ are positive constants depending only on $K, K'R^2$. Notice that we assert (16.92) for all $\rho \in (0, R)$ rather than $\rho \in (0, R/4)$ as in Theorem 16.15. We can do this because $|\nu| = 1$ (which means an inequality of the form (16.92) trivially holds for $\rho \in (R/4, R)$).

The estimate (16.92) can be used to obtain some rather strong regularity results for $\Stilde$. We first use (16.92) to deduce some facts about local non-parametric representations of $\Stilde$. Let $P$ be a linear isometry of $\R^3$ such that

\[
\vtilde(0) = P\nu(y, y_3) = (0, 0, 1)
\]

and let

\[
\Stilde = \{(\zeta, \zeta_3) \in \R^3 \mid (\zeta, \zeta_3) = P(x - y, x_3 - y_3),\ (x, x_3) \in \mathfrak{S}_{\theta R}^*(Y)\},
\]

where $\theta \in (0, 1)$. Since $\Stilde$ is a $C^2$ surface we of course know that for small enough $\theta$ there is a neighborhood $\mathcal U$ of $0 \in \R^2$ and a $C^2(\mathcal U)$ function $\tilde u$ with $D\tilde u(0) = 0$ and

\begin{equation}
\tag{16.93}
\Stilde = \operatorname{graph}\tilde u = \{(\zeta, \zeta_3) \in \R^3 \mid \zeta \in \mathcal U,\, \zeta_3 = \tilde u(\zeta)\}.
\end{equation}

Furthermore, writing

\[
\vtilde(\zeta) = (1 + |D\tilde u(\zeta)|^2)^{-1/2}(-D\tilde u(\zeta), 1), \qquad \zeta \in \mathcal U,
\]


% PDF page 439
\source{gilbarg-trudinger:page-0439}

\pagestyle{empty}



we have by (16.92) that

\[
|\vtilde(\zeta)-\vtilde(0)| \le C\theta^\alpha, \qquad \zeta\in\U,
\]

where $C,\alpha$ are as in (16.92). Consequently

\[
(1+|Du(\zeta)|^2)^{-1}|Du(\zeta)|^2+\bigl[(1+|Du(\zeta)|^2)^{-1/2}-1\bigr]^2\le (C\theta^\alpha)^2,\qquad \zeta\in\U.
\]

which implies

\[
(16.94)\qquad |Du(\zeta)|\le [1-(C\theta^\alpha)^2]^{-1/2}C\theta^\alpha<\frac12.
\]

provided $\theta$ is such that

\[
(16.95)\qquad C\theta^\alpha<\frac13.
\]

Because of (16.94), we can infer that a representation of the above type holds for any $\theta$ satisfying (16.95). This follows from the fact that if $\theta\in(0,1)$ is such that (16.93) and (16.95) hold, then, by using the smoothness of $\mathscr G$ together with (16.94), we can extend $\utilde$ so that a representation of the form (16.93) holds with $\mathscr G^{*}_{(\theta+\varepsilon)R}(Y)$ $(\varepsilon>0)$ in place of $\mathscr G^{*}_{\theta R}(Y)$. For later reference we also note that (16.94) implies

\[
(16.96)\qquad B_{\theta R/2}(0)\subset \U.
\]

We can now prove the non-trivial connectivity result of the following lemma.

\noindent\textbf{Lemma 16.17.}\ \emph{Suppose $G$ is $(K,K')$-quasiconformal. Then there is a constant $\theta\in(0,1)$, depending only on $K,K'R^2$, such that $\mathscr G_\rho(Y)$ is connected for each $\rho<\theta R$.}

\smallskip

\noindent\emph{Proof.}\ In the proof we will let $C_1,C_2,\ldots$ denote constants depending only on $K,K'R^2$. $\bar B_\sigma$, for $\sigma>0$, will denote the open ball $\{X\in\R^3\mid |X-Y|<\sigma\}$. Let $\theta\in(0,1)$ satisfy (16.95), $\rho=\theta R/2$, $\beta\in(0,\frac14)$ and define $\G_\beta$ to be the collection of those components of $\mathscr G_{\rho/2}(Y)$ which intersect the ball $\bar B_{\beta\rho}$. For each $\mathcal G\in\G_\beta$ we can find $Z\in\mathcal G\cap \bar B_{\rho/4}$ such that

\[
\mathcal G\subset \mathscr G_{\rho}^{*}(Z),
\]

and hence, replacing $Y$ by $Z$ and $R$ by $R/2$ in the discussion preceding the lemma, we see that $\mathcal G$ can be represented in the form (16.93), (16.94). Using such a non-parametric representation for each $\mathcal G\in\G_\beta$ and also using the fact that no two elements of $\G_\beta$ can intersect, it follows that the union of all the components $\mathcal G\in\G_\beta$ is contained in a region bounded between two parallel planes $\pi_1,\pi_2$ with

\[
(16.97)\qquad \dist(\pi_1,\pi_2)\le C_1(\beta+\theta^\alpha)\rho
\]

where $\alpha$ is as in (16.92).


% PDF page 440
\source{gilbarg-trudinger:page-0440}

\pagestyle{empty}



Our aim now is to show that, for suitable choices of $\beta$ and $\theta$ depending only
on $K$ and $K'R^2$, there is only one element (viz.\ $\mathscr{G}_{\rho/2}^*(Y)$) in $\Gcal_\beta$.
Suppose that in fact there are two distinct elements $\Gcal_1$, $\Gcal_2$ in $\Gcal_\beta$.
We can clearly choose $\Gcal_1$, $\Gcal_2$ to be adjacent in the sense that the volume
$\Vcal$ enclosed by $\Gcal_1$, $\Gcal_2$ and $\partial \Bbar_{\rho/2}$ contains no other
elements $\Gcal \in \Gcal_\beta$. Thus $\Vcal\cap \Bbar_{\beta\rho}$ consists entirely of points
above the graph $\G$ or entirely of points below $\G$; it is then evident that if the unit
normal $\mathbf{v}$ points out of (into) $\Vcal$ on $\Gcal_1$, then it also points out of (into)
$\Vcal$ on $\Gcal_2$. Furthermore, by (16.97) we have

\[
\text{volume }\Vcal = |\Vcal|\le C_2(\beta+\theta^\alpha)\rho^3,
\]

\begin{equation}
\tag{16.98}
\text{area }(\Vcal\cap \partial \Bbar_{\rho/2}) = A(\Vcal\cap \partial \Bbar_{\rho/2})
\le C_3(\beta+\theta^\alpha)\rho^2.
\end{equation}

An application of the divergence theorem over $\Vcal$ then gives

\[
\int_{\Gcal_1}\mathbf{v}\cdot\mathbf{v}\,dA+\int_{\Gcal_2}\mathbf{v}\cdot\mathbf{v}\,dA
= \pm\left\{\int_{\Vcal}\operatorname{div}\mathbf{v}\,dx\,dx_3
-\int_{\partial \Bbar_{\rho/2}\cap\Vcal}\eta\cdot\mathbf{v}\,dA\right\},
\]

where $\eta$ is the outward unit normal of $\partial \Bbar_{\rho/2}$. By (16.11) and (16.98) this gives

\[
A(\Gcal_1)+A(\Gcal_2)\le 2\int_{\Vcal}|H(x)|\,dx\,dx_3+C_3(\beta+\theta^\alpha)\rho^2.
\]

Also, by (16.98) and (16.91),

\[
\int_{\Vcal}|H(x)|\,dx\,dx_3
\le \left(\int_{\Vcal}H^2(x)\,dx\,dx_3\right)^{1/2}|\Vcal|^{1/2}
\]

\[
\le \left(\int_{\Bbar_{\rho/2}}H^2(x)\,dx\,dx_3\right)^{1/2}\{C_2(\beta+\theta^\alpha)\rho^3\}^{1/2}
\]

\[
\le (C_4\rho)^{1/2}\{C_2(\beta+\theta^\alpha)\rho^3\}^{1/2}
\]

\[
\le \sqrt{C_4C_2(\beta+\theta^\alpha)}\,\rho^2.
\]

Hence, provided $\beta+\theta^\alpha<1$, we have

\begin{equation}
\tag{16.99}
A(\Gcal_1)+A(\Gcal_2)\le C_5\sqrt{\beta+\theta^\alpha}\,\rho^2.
\end{equation}

On the other hand, by using a non-parametric representation as in (16.93),
(16.94) we infer that

\begin{equation}
\tag{16.100}
A(\Gcal)\ge C_6\rho^2
\end{equation}


% PDF page 441
\source{gilbarg-trudinger:page-0441}

\pagestyle{empty}



for each $\mathcal{G}\in\Gcal_\beta$, where $C_6>0$ is an absolute constant. (Since
$\mathscr{G}^*_{\rho/4}(Z)\subset\mathcal{G}$ we can deduce from (16.94) that (16.100)
holds with (for example) $C_6=\pi/64$. Inequalities (16.99) and (16.100) are clearly
contradictory if we choose $\beta$, $\theta$ small enough (but depending only on
$K$ and $K'R^2$). For such a choice of $\beta$, $\theta$ we thus have

\[
\mathscr{G}_{\beta\rho}(Y)=\mathscr{G}\cap\Bbar_{\beta\rho}
=\mathscr{G}_{\rho/2}(Y)\cap\Bbar_{\beta\rho}
=\mathscr{G}^*_{\rho/2}(Y)\cap\Bbar_{\beta\rho}.
\]

But by using a representation of the form (16.93), (16.94) for $\mathscr{G}^*_{\rho/2}(Y)$, we clearly
have $\mathscr{G}^*_{\rho/2}(Y)\cap\Bbar_{\beta\rho}$ connected. Thus $\mathscr{G}_{\beta\rho}(Y)=\mathscr{G}_{\theta R/2}(Y)$ is connected. The lemma
follows because $\beta$, $\theta$ depended only on $K$, $K'R^2$.
$\square$

Because of the above connectivity result we can now replace $\mathscr{G}^*_{\rho}(Y)$ in (16.92)
by $\mathscr{G}_{\rho}(Y)$ for $\rho\le \theta R$. However, since $|v|=1$, an inequality of the form (16.92)
is trivial for $\rho>\theta R$. Hence we have the following theorem.

\noindent\textbf{Theorem 16.18.}\ \emph{Suppose $G$ is $(K,K')$-quasiconformal. Then}

\begin{equation}
\tag{16.101}
\sup_{X\in\mathscr{G}_{\rho}(Y)}|v(X)-v(Y)|\le C(\rho/R)^\alpha,\qquad \rho\in(0,R),
\end{equation}

\emph{where $C$ and $\alpha$ are positive constants depending only on $K$, $K'R^2$.}

\smallskip

\noindent\emph{Remarks.}\ (i) The estimate (16.101) implies that

\begin{equation}
\tag{16.102}
|v(X)-v(\bar X)|\le 2^\alpha C(|X-\bar X|/R)^\alpha
\end{equation}

for all $X,\bar X\in\mathscr{G}_{R/4}(Y)$. This is seen by using (16.101) with $\bar X$ in place of $Y$ and with
$R/2$ in place of $R$.

(ii) If $K'=0$ and $\Omega=\R^2$, then we can let $R\to\infty$ in (16.101), thus showing
that $v(X)\equiv v(Y)$ on $\mathscr{G}$. That is, we have the following corollary.

\noindent\textbf{Corollary 16.19.}\ \emph{Suppose $G$ is $(K,0)$-quasiconformal and $\Omega=\R^2$. Then $u$ is a
linear function.}

Note that Corollary 16.19 can be directly deduced by letting $R\to\infty$ in (16.84),
without first proving (16.101) or even (16.92). However, we still require Lemma
16.16 to show that $\Lambda_1$ can be chosen to depend only on $K$.

\subsection*{16.7. Applications to Equations of Mean Curvature Type}

Here $\mathscr{G}$ will denote the graph of a solution $u$ of (16.63), (16.64), (16.65). The
remaining terminology will be as in Sections 16.4--16.6.

We first note that, since the Gauss map $G$ is automatically $(K,K')$-quasi-
conformal with $K$, $K'$ as in (16.89), we can immediately deduce Theorem 16.12 from
Corollary 16.19 above.


% PDF page 442
\source{gilbarg-trudinger:page-0442}

\setcounter{page}{430}

16. Equations of Mean Curvature Type

We next wish to show that Theorem 16.18 implies a bound for the principal
curvatures $\kappa_1$, $\kappa_2$ of $\mathcal{S}$, provided a suitable H\"older condition is imposed on the
coefficients $a^{ij}$, $b$. In order that this condition may be conveniently described, we
first extend the matrix $[a^{ij}]$ to be a $3\times 3$ matrix (cf. the procedure of Section 16.4)
by defining

\[
a^{3i}(x,z,p)=a^{i3}(x,z,p)=\sum_{j=1}^{2} p_j a^{ij}(x,z,p),\ i=1,2,
\qquad (x,z,p)\in\Omega\times\mathbb{R}\times\mathbb{R}^2,
\]

\[
a^{33}(x,z,p)=\sum_{i,j=1}^{2} p_i p_j a^{ij}(x,z,p).
\]

It will also be convenient to express $a^{ij}, b$ in terms of $X=(x,z)\in\Omega\times\mathbb{R}$ and
$v=(1+|p|^2)^{-1/2}(-p,1)$, by defining

\[
a_*^{ij}(X,v)=a^{ij}(x,z,p), \qquad i,j=1,2,3,
\]

\[
b_*(X,v)=(1+|p|^2)^{-1/2}b(x,z,p)
\]

for all $(x,z,p)\in\Omega\times\mathbb{R}\times\mathbb{R}^2$. Notice that these definitions give $a_*^{ij}, b_*$ on the sets
$(\Omega\times\mathbb{R})\times\{ \xi\in\mathbb{R}^3 \mid |\xi|=1,\ \xi_3>0\}$. (In the case when equation (16.63) arises as the
non-parametric Euler equation of an elliptic parametric variational problem, we
show in Appendix 16.8 that the functions $a_*^{ij}, b_*$ arise quite naturally.) We now
assume that the functions $a_*^{ij}, b_*$ satisfy the H\"older conditions

\begin{equation}
\begin{aligned}
\left|a_*^{ij}(X,v)-a_*^{ij}(\bar X,\bar v)\right| &\le \mu_1\{(|X-\bar X|/R)^\beta+|v-\bar v|^\beta\},\\
\left|b_*(X,v)-b_*(\bar X,\bar v)\right| &\le \mu_2\{(|X-\bar X|/R)^\beta+|v-\bar v|^\beta\},
\end{aligned}
\tag{16.103}
\end{equation}

for all $X,\bar X\in\Omega\times\mathbb{R}$ and all $v,\bar v\in\{\xi\in\mathbb{R}^3\mid |\xi|=1,\ \xi_3>0\}$, where $\mu_1,\mu_2$ and $\beta$ are
constants with $\beta\in(0,1)$.

\textbf{Theorem 16.20.} \textit{Suppose (16.63), (16.64), (16.65) and (16.103) hold. Then if $\kappa_1,\ \kappa_2$
are the principal curvatures of $\mathcal{S}$ at $Y$, we have}

\begin{equation}
\kappa_1^2+\kappa_2^2\le C/R^2,
\tag{16.104}
\end{equation}

\noindent\textit{where}\qquad $C=C(\gamma,\mu R,\mu_1,\mu_2 R,\beta)$.

\textit{Proof.} Choose $\theta$ small enough to ensure that $\mathcal{S}_{\theta R}(Y)$ is connected (Lemma 16.17)
and can be represented in the form (16.93) with (16.95) (and hence (16.94)) holding.
By applying the discussion of Section 16.4 to both $\mathcal{S}$ and the transformed surface
$\tilde{\mathcal{S}}$ of (16.93), it then follows that $\tilde u$ satisfies an equation of the form

\begin{equation}
\tilde a^{ij}(\zeta)D_{ij}\tilde u+\tilde b(\zeta)=0
\tag{16.105}
\end{equation}


% PDF page 443
\source{gilbarg-trudinger:page-0443}

on $\mathcal U$, where

\[
|\xi|^2-\frac{(D\tilde u\cdot \xi)^2}{1+|D\tilde u|^2}
\leq \tilde a^{ij}\xi_i\xi_j
\leq \gamma\left[|\xi|^2-\frac{(D\tilde u\cdot \xi)^2}{1+|D\tilde u|^2}\right]
\]

for all $\xi\in \mathbb R^2$, and

\[
|\tilde b|\leq \mu\sqrt{1+|D\tilde u|^2}.
\]

Consequently, by (16.94), we have for all $\xi\in B=B_{R/2}(0)$

\begin{equation}\tag{16.106}
|\xi|^2/2\leq \tilde a^{ij}(\zeta)\xi_i\xi_j\leq \gamma|\xi|^2,
\qquad
|\tilde b(\zeta)|\leq 2\mu.
\end{equation}

By (16.103), (16.102) and the discussion in Section 16.4, it also follows that we can
assume the H\"older estimates

\begin{equation}\tag{16.107}
\left|\tilde a^{ij}(\zeta)-\tilde a^{ij}(\bar\zeta)\right|
\leq C(|\zeta-\bar\zeta|/R)^{x\beta},
\qquad \zeta,\bar\zeta\in B,\ i,j=1,2
\end{equation}
\[
\left|\tilde b(\zeta)-\tilde b(\bar\zeta)\right|
\leq C(|\zeta-\bar\zeta|/R)^{x\beta},
\qquad \zeta,\bar\zeta\in B.
\]

Furthermore it is clear from (16.94), and the fact that $\tilde u(0)=0$, that

\begin{equation}\tag{16.108}
\sup_B |\tilde u|\leq R.
\end{equation}

Then, by using Schauder's interior estimate (Theorem 6.2) in conjunction with
(16.106), (16.107) and (16.108), we infer that

\[
|\tilde u|_{2,x;B}^*\leq C\{R+\mu R^2+R^{2+x\beta}\mu_2R^{-x\beta}\}
\]
\[
\leq CR,
\]

where $C=C(\gamma,\mu R,\mu_1,\mu_2R,\beta)$. In particular we have

\[
\sum_{i,j=1}^2 |D_{ij}\tilde u(0)|\leq C/R,
\]

whence the theorem follows by (16.87). $\square$

The H\"older estimate (16.101) can also be used to obtain gradient estimates for
solutions $u$ of (16.63), (16.64), (16.65). The following theorem deals with the homo-
geneous case, when no smoothness or continuity restrictions will be imposed on
the coefficients.


% PDF page 444
\source{gilbarg-trudinger:page-0444}

\pagestyle{empty}


\setcounter{section}{16}

\textbf{Theorem 16.21.} \emph{Suppose (16.63), (16.64) hold, $b \equiv 0$, and suppose the functions
$a^{ij}(x,u,Du)$, $i,j=1,2$ are measurable on $\Omega$. Then}

\begin{equation}
|Du(y)| \le C_1 \exp(C_2 m_R / R).
\tag{16.109}
\end{equation}

\emph{where}

\[
m_R = \sup_{B_R(y)} (u-u(y))
\]

\emph{and $C_1 = C_1(\gamma)$, $C_2 = C_2(\gamma)$.}

\emph{Proof.} As in Theorem 16.20 we assume $\theta$ is small enough to ensure that
$\Xi_{\theta R}(Y)$ is connected and that the representation (16.93), (16.94) holds.
Notice that since $b \equiv 0$ we can choose $\theta$ to depend only on $\gamma$.
Also, we can here suppose that the matrix of $P$, $[p_{ij}]$, is chosen such that
$p_{32}=0$. Hence we have

\begin{equation}
(1+|Du(x)|^2)^{-1/2} = (1+|D\tilde u(\zeta)|^2)^{-1/2}(p_{31}D_1\tilde u(\zeta)-p_{33}).
\tag{16.110}
\end{equation}

Now (since $b \equiv 0$), equation (16.105) can be written in the form

\[
\alpha D_1 \tilde u + 2\beta D_{12}\tilde u + D_{22}\tilde u = 0
\]

\emph{where} $\alpha = \tilde a_{11}/\tilde a_{22}$ \emph{and} $\beta = \tilde a_{12}/\tilde a_{22}$.
We now multiply each side of this equation by $p_{31}D_1\varphi$, where
$\varphi \in C_0^2(\mathcal U)$ and integrate over $\mathcal U$.
Making use of the relations

\[
\int_{\mathcal U} D_{22}\tilde u\,D_1\varphi\,d\zeta
= -\int_{\mathcal U} D_2\tilde u\,D_{12}\varphi\,d\zeta
= \int_{\mathcal U} D_{21}\tilde u\,D_2\varphi\,d\zeta
\]

\emph{and writing}

\[
\psi = p_{31}D_1\tilde u - p_{33},
\]

\emph{we then obtain}

\[
\int_{\mathcal U} \left(\alpha D_1\psi\,D_1\varphi + 2\beta D_2\psi\,D_1\varphi + D_2\psi\,D_2\varphi\right)\,d\zeta = 0.
\]

\emph{That is,} $\psi$ \emph{is a weak solution of the uniformly elliptic equation}

\[
D_1(\alpha D_1\psi + 2\beta D_2\psi) + D_2(D_2\psi)=0
\qquad \text{on } \mathcal U.
\]

\emph{Furthermore, we have $\psi > 0$ on $\mathcal U$ by (16.110). Hence we can apply the Harnack
inequality, Theorem 8.20, to $\psi$, thus giving}

\begin{equation}
\sup_{B_{R/2}(0)} \psi \le C \inf_{B_{R/2}(0)} \psi.
\tag{16.111}
\end{equation}


% PDF page 445
\source{gilbarg-trudinger:page-0445}

where $C=C(\gamma)$. But because of (16.110) and (16.94) we know that

\[
(1+|Du(x)|^2)^{-1/2}\leq \psi(\zeta)\leq 2(1+|Du(x)|^2)^{-1/2}
\]

on $\mathcal{U}$, and hence defining $v$ on $\Omega\times\mathbb{R}$ by

\[
v(x,z)=(1+|Du(x)|^2)^{1/2}, \qquad x\in \Omega, z\in \mathbb{R},
\]

we can deduce from (16.111) that

\[
\sup_{X\in G_{R/2}} v(X)\leq C \inf_{X\in G_{R/2}} v(X),
\]

where $C=C(\gamma)$. By varying $Y$ it clearly follows that there is a number $\chi\in (0,\frac12)$,
depending only on $\gamma$, such that

\[
\text{(16.112)} \qquad v(X_1)\leq Cv(X_2)
\]

whenever $X_1=(x^{(1)},u(x^{(1)}))$ and $X_2=(x^{(2)},u(x^{(2)}))$ are such that $|X_1-X_2|\leq \chi R$
and $x^{(1)}, x^{(2)}\in B_{R/2}(y)$. Now let

\[
B_{\chi R}^{+}(y)=\{x\in B_{\chi R}\mid u(x)>u(y)\}
\]

and let $Y_1=(y^{(1)},u(y^{(1)}))$ and $Y_2=(y^{(2)},u(y^{(2)}))$ be such that $y^{(1)}, y^{(2)}\in \overline{B_{\chi R}^{+}(y)}$
and

\[
|Du(y^{(1)})|=\sup_{\overline{B_{\chi R}^{+}(y)}} |Du|,
\]

\[
|Du(y^{(2)})|=\inf_{\overline{B_{\chi R}^{+}(y)}} |Du|.
\]

Take a sequence $X_1,\ldots,X_N$ of points in $\mathcal{G}\cap (\overline{B_{\chi R}^{+}(y)}\times \mathbb{R})$ such that
$|X_{i+1}-X_i|<\chi R$, $i=1,\ldots,N-1$, and such that $X_1=Y_1$, $X_N=Y_2$. Clearly, repeated application
of (16.111) implies

\[
v(Y_1)\leq C^N v(Y_2);
\]

that is

\[
\sup_{\overline{B_{\chi R}^{+}(y)}} \sqrt{1+|Du|^2}\leq C^N \inf_{\overline{B_{\chi R}^{+}(y)}} \sqrt{1+|Du|^2}.
\]

However, it is clear that we can choose $N$ so that

\[
N\leq C(m_R/R+1),
\]

% PDF page 446
\source{gilbarg-trudinger:page-0446}

where $C=C(\gamma)$. Hence we obtain
\[
\sup_{B^+_{x_R}(y)} \sqrt{1+|Du|^2}\leq \{C_1 \exp(C_2 m_R/R)\} \inf_{B^+_{x_R}(y)} \sqrt{1+|Du|^2},
\]
where $C_1=C_1(\gamma)$, $C_2=C_2(\gamma)$. Finally, Theorem 16.21 follows by using the fact that
\[
\inf_{B^+_{x_R}(y)} |Du|\leq \chi^{-1} m_R/R.
\]
(See problem 16.5.) $\square$

The H\"older estimate for $v$ can also be used in the non-homogeneous case to deduce a gradient estimate for $u$, but in this case it is necessary to impose Lipschitz restrictions on the functions $a_*^{ij}$, $b_*$, introduced above. The interested reader is referred to [S1 4].

\bigskip

\noindent 16.8. \textbf{Appendix: Elliptic Parametric Functionals}

Let $\Omega$ be a bounded domain in $\mathbb{R}^2$ and consider the functional $I$, defined for $C^1$ mappings $\mathbf{Y}=(Y_1,\,Y_2,\,Y_3): \overline{\Omega}\to \mathbb{R}^3$ by
\[
I(\mathbf{Y})=\int_\Omega G(x,\mathbf{Y},D_1\mathbf{Y},D_2\mathbf{Y})\,dx,
\tag{16.113}
\]
where $G=G(x,X,p)$ is a given continuous function of $(x,X,p)\in \mathbb{R}^2\times \mathbb{R}^3\times \mathbb{R}^6$.
(Here of course $D_i\mathbf{Y}=(D_iY_1,\,D_iY_2,\,D_iY_3)$ for $i=1,2$). Now let us consider the
possibility that $I$ remains invariant under orientation preserving diffeomorphisms
of $\mathbb{R}^2$; that is, whenever $\psi$ is a diffeomorphism of $\mathbb{R}^2$ onto itself with positive
Jacobian, we would have
\[
\int_{\Omega'} G(\zeta,\widetilde{\mathbf{Y}}(\zeta),D_1\widetilde{\mathbf{Y}}(\zeta),D_2\widetilde{\mathbf{Y}}(\zeta))\,d\zeta
=\int_\Omega G(x,\mathbf{Y}(x),D_1\mathbf{Y}(x),D_2\mathbf{Y}(x))\,dx,
\]
where $\Omega'=\psi(\Omega)$ and $\widetilde{\mathbf{Y}}=\mathbf{Y}\circ \psi^{-1}$. A simple computation (cf. [MY 5], p.\ 349)
shows that this would be true for all such diffeomorphisms $\psi$ and domains $\Omega$ if
and only if there is a real-valued function $F$ on $\mathbb{R}^3\times \mathbb{R}^3$ such that
\[
G(x,X,p)=F(X,P), \qquad (x,X,p)\in \mathbb{R}^2\times \mathbb{R}^3\times \mathbb{R}^6,
\tag{16.114}
\]
where
\[
P=(p_3p_5-p_2p_6,\; p_1p_6-p_3p_4,\; p_2p_4-p_1p_5),
\]


% PDF page 447
\source{gilbarg-trudinger:page-0447}

\pagestyle{empty}



\noindent\textbf{16.8 Appendix: Elliptic Parametric Functionals}

\medskip

and

\[
(16.115)\qquad F(X,\lambda q)=\lambda F(X,q), \qquad (X,q)\in\R^3\times\R^3,\qquad \lambda>0.
\]

Note in particular that (16.114) implies that $G(x,X,p)$ cannot depend on $x$; that
is, $G(x,X,p)=G(0,X,p)$ for $(x,X,p)\in\R^2\times\R^3\times\R^6$. In case $p=(D_1\Y,D_2\Y)$,
where $\Y$ is a $C^1$ map from $\overline{\Omega}$ into $\R^3$, $P$ is given by

\[
P=(D_1Y_3\cdot D_2Y_2-D_1Y_2\cdot D_2Y_3,\; D_1Y_1\cdot D_2Y_3-D_1Y_3\cdot D_2Y_1,
\]
\[
\hphantom{P=(}\; D_1Y_2\cdot D_2Y_1-D_1Y_1\cdot D_2Y_2).
\]

As is well known, in case $\Y$ is one-to-one and such that the Jacobian matrix
$[D_jY_i(x)]$ has rank 2 for each $x\in\Omega$, this last identity can be written

\[
P=\chi\nu,
\]

where $\nu$ is the unit normal of the embedded surface $\G=\{Y(x)\mid x\in\Omega\}$ and $\chi$ is
the area magnification factor of the mapping $\Y$. Thus, assuming that we orient $\G$
with unit normal $\nu$ such that $\chi>0$, we can write

\[
I(\Y)=\int_{\G} F(X,\nu(X))\,dA(X);
\]

that is, we can express $I(\Y)$ completely in terms of the oriented surface $\G$ and
independently of the particular mapping $\Y$ that is used to represent $\G$. Through
this discussion we are led to consider the functional $J$, defined for any smooth
oriented surface $\G$ in $\R^3$ having finite area by

\[
(16.116)\qquad J(\G)=\int_{\G} F(X,\nu(X))\,dA(X);
\]

this functional has the property that $J(\G)=I(\Y)$ whenever $\Y$ is a one-to-one $C^1$
mapping from $\Omega$ into $\R^3$ such that $[D_jY_i]$ has rank 2 at each point of $\Omega$ and
$\G=\{Y(x)\mid x\in\Omega\}$.

If $F$ satisfies (16.115), we call a functional of the form (16.116) a \emph{parametric}
\emph{functional}. The functional $J$ is called elliptic if $F$ is $C^2$ on $\R^3\times(\R^3-\{0\})$ and if
the convexity condition

\[
(16.117)\qquad D_{q_i q_j}F(X,q)\,\xi_i\xi_j\ge |q|^{-1}|\xi'|^2,\qquad
\xi'=\xi-\left(\xi\cdot\frac{q}{|q|}\right)\frac{q}{|q|}
\]

holds for all $X\in\R^3$, $q\in\R^3-\{0\}$ and $\xi\in\R^3$. Notice that, up to a scalar factor,
(16.117) is the strongest convexity condition possible for $F$ in view of the homo-
geneity condition (16.115).


% PDF page 448
\source{gilbarg-trudinger:page-0448}

If we now consider a non-parametric surface $\mathfrak{S}$ given by

\[
\mathfrak{S}=\{(x,u(x))\in \mathbb{R}^{3}\mid x\in \Omega\},
\]

where $u\in C^{2}(\bar{\Omega})$, then, taking $\nu$ to be the downward unit normal
\[
(Du,-1)/\sqrt{1+|Du|^{2}},
\]
we have

\[
J(\mathfrak{S})=\int_{\Omega}F(x,u(x),Du(x),-1)\,dx.
\]

Notice that here we have used the relation $dA=\sqrt{1+|Du|^{2}}\,dx$. The expression on the
right can be considered as a non-parametric functional, defined for any $u\in C^{2}(\bar{\Omega})$.
The Euler-Lagrange equation for this non-parametric functional is

\[
\sum_{i=1}^{2}D_{i}[D_{q_{i}}F(x,u,Du,-1)]-D_{x_{3}}F(x,u,Du,-1)=0.
\]

By using the chain rule and the homogeneity condition (16.115), one can easily check that this
equation can be written in the form

\[
Qu=a^{ij}(x,u,Du)D_{i,j}u+b(x,u,Du)=0,
\]

where

\[
(16.118)\qquad a^{ij}(x,u,Du)=\nu D_{q_{i}q_{j}}F(x,u,Du,-1),\qquad i,j=1,2,
\]

\[
(16.119)\qquad b(x,u,Du)=\nu\sum_{i=1}^{3}D_{q_{i}x_{i}}F(x,u,Du,-1),
\]

with $\nu=\sqrt{1+|Du|^{2}}$. By using (16.115), (16.117) it is not difficult to check that
(16.64), (16.65) hold with constants $\gamma$ and $\mu$ depending on $F$. That is, the non-
parametric Euler-Lagrange equation for an elliptic parametric functional is an equation of mean
curvature type.

Finally we wish to point out that the functions $a^{ij}_{*}, b_{*}$ introduced in Section 16.7 have a
natural interpretation in the present context. In fact one can easily check that in case $a^{ij}, b$
are as in (16.118), (16.119), then $a^{ij}_{*}, b_{*}$ are given by

\[
a^{ij}_{*}(X,\nu)=D_{q_{i}q_{j}}F(X,\nu),\qquad i,j=1,2,3.
\]

\[
b_{*}(X,\nu)=\sum_{i=1}^{3}D_{q_{i}x_{i}}F(X,\nu);
\]

furthermore the conditions (16.103) hold automatically with $\mu_{1}, \mu_{2}$ determined
by $F$, provided that $F\in C^{3}(\mathbb{R}^{3}\times (\mathbb{R}^{3}-\{0\}))$.

% PDF page 449
\source{gilbarg-trudinger:page-0449}

\begin{center}
\textbf{Notes}
\end{center}

\textbf{Notes}

\medskip

The interior gradient bound, Theorem 16.5, for the minimal surface equation was
discovered, in the case of two variables, by Finn [FN 2] and in the general case by
Bombieri, De Giorgi and Miranda [BDM]. An interior gradient bound for the
general prescribed mean curvature equation (16.1) was first established by Lady-
zhenskaya and Ural'tseva [LU 6]. The methods of both papers [BDM] and [LU 6]
depended upon an isoperimetric inequality of Federer and Fleming (see [FE])
and a resulting Sobolev inequality (see [MD 1] and [LU 6]). Our derivation of
Theorem 16.5, together with the relevant preparatory material in Section 16.1, is
adapted from the work of Michael and Simon [MS] and Trudinger [TR 6, 8].
The key idea of employing methods analogous to classical potential theory was
invoked in an (unpublished) simplification by Michael of the Sobolev inequality
in [MD 1] and subsequently utilized in [MS] and [TR 6, 8]. The integration by
parts formula, Lemma 16.1, is due to Morrey [MY 3]. The mean value inequality,
Lemma 16.2, is due to Michael and Simon [MS] and its extension, Lemma 16.3,
is essentially given in [TR 8]. The extension of Morrey's estimate to hypersurfaces,
Lemma 16.4, is given in [SI 4]. The process in Section 16.2, by which we conclude
an interior gradient bound from the potential type inequality, Lemma 16.3, is
taken from [TR 6, 8].

The existence theorems 16.9 and 16.11, for the cases of smooth boundary data,
were established by Serrin [SE 4]. Theorem 16.10 essentially appears in Giaquinta
[GI]. Note that the results of Section 16.3 can be approached by the variational
method described in Section 10.5. In recent years the associated variational
problems have also been studied in the space of functions of bounded variation,
(see for example [BG 2], [GE], [GI], [MD 5]). A further approach to the equation
of prescribed mean curvature is presented in [TE].

The treatment of equations of mean curvature type in Section 16.4 is based on
[SI 4, 6]. The pioneering work on two dimensional equations of mean curvature
type was done by Finn [FN 1, 2] who treated the case $a^{ij}(x, z, p)\equiv a^{ij}(p)$ and
$b\equiv 0$. Finn called his equations ``equations of minimal surface type'' and stated
the structure conditions for the coefficient matrix somewhat differently (but
equivalently) to (16.64). The first gradient estimates were obtained in [FN 1, 2].
For a more specialized class of equations than those considered in [FN 2], (in
fact for equations of the form (16.63) with $b\equiv 0$ and $a^{ij}$ as in (16.118) for some $F$
with $F(X, p)\equiv F(0, p)$), refinements (including an inequality like that of Theorem
16.21) were obtained by Jenkins and Serrin [JS 1]. The inequality (16.112) seems
to be new. A curvature estimate like that in Theorem 16.20 was originally obtained
for the minimal surface equation by Heinz [HE 1] and strengthened by E. Hopf
[HO 4] and R. Osserman [OM 1]. Jenkins [JE] and Jenkins and Serrin [JS 1]
obtained similar curvature estimates for equations of the form (16.63) in the case
when $b\equiv 0$ and when $a^{ij}$ had the form (16.118) for some $F$ with $F(X, p)\equiv F(0, p)$.
The estimates in [HO 4], [OR 1], [JE] and [JS 1] were in fact obtained in the
stronger form

\begin{equation}
(16.120)\qquad (1+|\Du(y)|^2)(\kappa_1^2+\kappa_2^2)(y)\leq C/R^2.
\end{equation}


% PDF page 450
\source{gilbarg-trudinger:page-0450}


For the special class of equations dealt with in [JE] and [JS 1], the methods of
Sections 4 to 7 can be modified to also give an inequality of the form (16.120), as
shown in [SI 4]. In the case $b\neq 0$ it can easily be shown that an estimate like
(16.120) cannot hold in general. The only curvature estimate for the case $b\neq 0$ prior
to those obtained here and in [SI 4], as far as the authors are aware, was the result
of Spruck [SP] for the constant mean curvature equation. It should also be
mentioned that results for special classes of parametric surfaces have been obtained
in [OM 1], [JE], [JS 1] and [SI 6].

The general Bernstein type result of Corollary 16.19 settles a question raised
by Osserman [OM 2, p. 137]; such a result is well known (and there are many
proofs) in the case of the minimal surface equation; (see for example [NT]). For
the class of equations dealt with in [JE], Jenkins obtained such a result by letting
$R\to\infty$ in (16.120). The question of whether the Bernstein theorem for the
minimal surface equation carried over to $\mathbb{R}^n$, $n>2$, provided a great impetus to the
study of higher dimensional minimal surfaces. It was eventually shown to be true
in case $n\le 7$ by Simons [SM], who used some ideas of Fleming [FL] and De Giorgi
[DG 2]. It was shown to be false in case $n>7$ by Bombieri, De Giorgi and Giusti
[BDG]. Curvature estimates of the type established in Theorem 16.20 were shown
also to hold for the minimal surface equation when $n\le 7$ by L. Simon [SI 5], and
these imply Bernstein's theorem for $n\le 7$.

For a thorough account of two dimensional minimal surfaces the reader is
referred to the book [NT]. For recent developments in the higher dimensional
theory, see [SI 7, 8], [GT 5].

\bigskip
\noindent\textbf{Problems}

\medskip
\noindent\textbf{16.1.}\quad Show that by replacing $n$ by some number $\alpha<n$ in the choice of the function
$\chi$ in the derivation of Lemma 16.2 we obtain, instead of (16.27) and (16.40),
the relations:

\begin{align*}
&\left(1-\frac{\alpha}{n}\right)\int_{\mathcal S_R} r^{-\alpha} g\, dA+\frac{\alpha}{n}\int_{\mathcal S_R} r^{-\alpha}(1-|\delta r|^2)g\, dA\tag{16.121}\\
&\qquad = R^{-\alpha}\int_{\mathcal S_R} g\, dA+\int_{\mathcal S_R}(r^{-\alpha}-R^{-\alpha})gH\nu\cdot (x-y)\, dA\\
&\qquad\quad -\frac{1}{n}\int_{\mathcal S_R}(r^{-\alpha}-R^{-\alpha})(x-y)\cdot \delta g\, dA
\end{align*}

for $g\in C^1(\mathcal U)$;


% PDF page 451
\source{gilbarg-trudinger:page-0451}

\begin{quote}
\setcounter{equation}{122}
\begin{equation}\tag{16.122}
\left(1-\frac{\alpha}{n}\right)\int_{\mathcal S} r^{-\alpha}g\,dA+\frac{\alpha}{n}\int_{\mathcal S}r^{-\alpha}(1-|\delta r|^2)g\,dA
=\int_{\mathcal S}r^{-\alpha}gH\,\nabla\!\cdot(x-y)\,dA-\frac{1}{n}\int_{\mathcal S}r^{-\alpha}(x-y)\cdot\delta g\,dA
\end{equation}

for $g\in C^1(\mathcal U)$.

\medskip

\textbf{16.2.} Using (16.40) and (16.121) derive the following Sobolev inequalities for
$g\in C_0^1(\mathcal U)$:

\begin{equation}\tag{16.123}
\left[\int_{\mathcal S}|g|^p\,dA\right]^{1/p}\le C(n,p)(\operatorname{diam}\mathcal U)^{1-n/p}\,[A(\mathcal S)]^{1/p}\int_{\mathcal S}|\delta g|\,dA
\end{equation}

for $p<n/(n-1)$, $H\equiv0$, where

\[
C(n,p)=\frac{1}{n\omega_n}\left[\frac{n}{n-(n-1)p}\right]^{1/p}.
\]

\begin{equation}\tag{16.124}
\left[\int_{\mathcal S}|g|^2\,dA\right]^{1/2}\le \frac{1}{\sqrt{\pi}}\int_{\mathcal S}(|\delta g|+|Hg|)\,dA
\end{equation}

for $n=2$. (Note that (16.123) is established for $p=n/(n-1)$ in [MSI].)

\medskip

\textbf{16.3.} Let $g$ be a non-negative subharmonic function on a $C^2$ hypersurface
$\mathcal S\subset\mathbb R^{n+1}$. Derive the following generalization of the mean value inequality (16.21):

\[
g(y)\le
\begin{cases}
\dfrac{1}{\pi R^2}\int_{\mathcal S_R}g\,dA+\dfrac{1}{4\pi}\int_{\mathcal S_R}gH^2\,dA & \text{if } n=2,\\[1.2em]
\dfrac{1+C(n)\bigl[H_0R+(H_0R)^n\bigr]}{\omega_n R^n}\int_{\mathcal S_R}g\,dA & \text{if } n>2,\, H_0=\sup_{\mathcal S}|H|.
\end{cases}
\]

\medskip

\textbf{16.4} Using Corollary (16.6) and the Harnack inequality, Theorem 8.28, show
that a solution of the minimal surface equation in $\mathbb R^n$ which is bounded above by a
linear function must itself be linear.

\medskip

\textbf{16.5.} Derive the inequality

\[
\inf_{B^{+}_{\chi R}(y)}|Du|\le \chi^{-1}m_R/R
\]

used in the proof of Theorem 16.21.
\end{quote}

% PDF page 452
\source{gilbarg-trudinger:page-0452}


\noindent\textbf{16.6.}\quad Using (16.60) show that the constants $C_1$ and $C_2$ in Theorem 16.5 and
Corollary 16.6 need only depend on $n$ and $d^2 H_1$.

\medskip

\noindent\textbf{16.7.}\quad By adapting the Perron process described in Chapter 2, show that condition
(16.60) is sufficient to guarantee the existence of a classical solution of the prescribed
mean curvature equation, $\mathcal M u = H(1+|Du|^2)^{3/2}$, in a domain $\Omega$. Show that
condition (16.60) with $\varepsilon_0 = 0$ is necessary for the existence of a classical solution.
It turns out that this condition is also sufficient [GT 3].


% PDF page 453
\source{gilbarg-trudinger:page-0453}


\noindent Chapter 17

\bigskip

\noindent {\Large Fully Nonlinear Equations}

\vspace{2.5cm}

\noindent In this chapter we consider the solvability of the classical Dirichlet problem for
certain types of fully nonlinear elliptic equations; that is, nonlinear elliptic equations
that are not quasilinear. A general second-order equation, on a domain $\Omega$ in $\mathbb{R}^n$,
can be written in the form,

\medskip

\noindent (17.1)\quad $F[u] = F(x,u,Du,D^2u) = 0,$

\medskip

\noindent where $F$ is a real function on the set $\Gamma = \Omega \times \mathbb{R} \times \mathbb{R}^n \times \mathbb{R}^{n\times n}$, where $\mathbb{R}^{n\times n}$ denotes
the $n(n+1)/2$ dimensional space of real symmetric $n \times n$ matrices. We denote
points in $\Gamma$ typically by $\gamma = (x, z, p, r)$ where $x \in \Omega$, $z \in \mathbb{R}$, $p \in \mathbb{R}^n$ and $r \in \mathbb{R}^{n\times n}$. When
$F$ is an affine function of the $r$ variables, the equation (17.1) is called quasilinear;
otherwise, it is called fully nonlinear. When $F$ is differentiable with respect to the $r$
variables, the following definitions extend those in Chapter 10:

\begin{quote}
The operator $F$ is elliptic in a subset $\mathcal U$ of $\Gamma$ if the matrix $[F_{ij}(\gamma)]$, given by
\end{quote}

\[
F_{ij}(\gamma) = \frac{\partial F}{\partial r_{ij}}(\gamma), \qquad i, j = 1, \ldots, n,
\]

\noindent is positive for all $\gamma = (x, z, p, r) \in \mathcal U$. Letting $\lambda(\gamma)$, $\Lambda(\gamma)$ denote, respectively, the
minimum and maximum eigenvalues of $[F_{ij}(\gamma)]$, we call $F$ uniformly elliptic (strictly
elliptic) in $\mathcal U$, if $\Lambda/\lambda\,(1/\lambda)$ is bounded in $\mathcal U$. If $F$ is elliptic (uniformly elliptic, strictly
elliptic) in the whole set $\Gamma$, then we simply say that $F$ is elliptic (uniformly elliptic,
strictly elliptic) in $\Omega$. If $u \in C^2(\Omega)$ and $F$ is elliptic (uniformly elliptic, strictly elliptic)
on the range of the mapping $x \mapsto (x,u(x),Du(x),D^2u(x))$, we say that $F$ is elliptic
(uniformly elliptic, strictly elliptic) with respect to $u$.

\bigskip

\noindent \textsc{Examples}

\medskip

\noindent (i) \textit{The Monge-Amp\`ere Equation:}

\medskip

\noindent (17.2)\quad $F[u] = \det\, D^2u - f(x) = 0.$
